\documentclass[12pt, a4paper]{article}
\usepackage[T1]{fontenc}
\usepackage{lmodern}
\usepackage{amsmath, amsthm, amssymb, mathtools}
\usepackage[margin=1in]{geometry}
\usepackage{xr-hyper}
\usepackage[colorlinks=true,linkcolor=blue,citecolor=blue]{hyperref}
\usepackage{cleveref}
\usepackage{graphicx, subcaption}
\usepackage{natbib}
\usepackage{booktabs}
\usepackage{enumerate}

% Code references: scaled monospace with line-breaking
\newcommand{\code}[1]{\texttt{\small #1}}

\newtheorem{theorem}{Theorem}[section]
\newtheorem{lemma}[theorem]{Lemma}
\newtheorem{proposition}[theorem]{Proposition}
\newtheorem{corollary}[theorem]{Corollary}
\newtheorem{definition}{Definition}[section]
\newtheorem{remark}{Remark}[section]
\newtheorem{conjecture}[theorem]{Conjecture}
\newtheorem{derivation}[theorem]{Derivation}
\newtheorem{observation}[theorem]{Observation}
\numberwithin{equation}{section}

\newcommand{\dd}{\mathrm{d}}
\newcommand{\seriestitle}{Why Rotating Fluids Make Polygons}


\externaldocument[I-]{../paper-1-mathematics/main}
\externaldocument[II-]{../paper-2-physics/main}
\externaldocument[IV-]{../paper-4-field-theory/main}
\externaldocument[V-]{../paper-5-cosmology/main}
\externaldocument[VI-]{../paper-6-discussion/main}

\title{\seriestitle\\[6pt]\large Paper III: Gravitational Theory}
\author{Gordon Speagle}
\date{}

\begin{document}
\maketitle

\begin{abstract}
Building on the vortex--gravity isomorphism established in Paper~II,
we show that in $2{+}1$D the Havelock eigenvalue
structure, combined with the Onsager selection principle
at negative temperature, implies the vacuum Einstein equations
$R_{\mu\nu} = \Lambda g_{\mu\nu}$.
The argument rests on three conditions: H1~(the $\csc^2$
uniqueness, a theorem), H2~(strict concavity of the polygon
entropy functional with Onsager selection of the unique
maximum), and H3~(Hadamard universality, a theorem);
the polygon--BTZ phase transition, the orbifold--Havelock
correspondence at $c = 12\,b(N)$, and the Wheeler--DeWitt
equation on the palindromic minisuperspace provide the
supporting infrastructure.
In $d \ge 3$ the characterisation extends, providing a new
criterion for Einstein manifolds.
\end{abstract}

%% ====================================================================
\section{Introduction and summary}
\label{sec:intro-gravity}
%% ====================================================================

Papers~I and~II develop the Havelock decomposition as
pure mathematics and fluid dynamics.
This paper connects the eigenvalue
structure to gravity.
The starting point is the observation (Paper~II) that
$N$ equal point masses on a spatial slice of
$2{+}1$D gravity interact via the same Green's function
as point vortices
(Proposition~\ref{II-prop:grav-stability}), giving identical Havelock
eigenvalues $\lambda_m = C_1(\xi) - m(N{-}m)/2$.
The conformal invariance of the flat-plane backreaction
(Proposition~\ref{II-prop:conformal-invariance}) ensures that
$N_{\mathrm{crit}} = 7$ is independent of the deficit angle, a
simplification that makes the algebraic-number thresholds of
Paper~I into gravitational phase transitions.

This paper develops the consequences:
\begin{itemize}
  \item \S\ref{sec:graviton}--\S\ref{sec:spectral-bound}:
    the graviton at $N = 7$ and the spectral bound as a
    quantum gravitational bound.
  \item \S\ref{sec:holographic}--\S\ref{sec:ope-identification}:
    holographic interpretation and the conformal block identification.
  \item \S\ref{sec:btz}:
    the polygon--BTZ phase transition, entropy accounting,
    and Lorentzian continuation.
  \item \S\ref{sec:wdw}--\S\ref{sec:three-layer}:
    the Wheeler--DeWitt equation and the three-layer decomposition.
  \item \S\ref{sec:cft3}--\S\ref{sec:triangle}:
    higher-dimensional extension and triangle universality.
  \item \S\ref{sec:gravity-partition}--\S\ref{sec:lichnerowicz}:
    the gravity partition function and the derivation of the
    Einstein equations.
\end{itemize}

\paragraph{Main result.}
The central theorem of this paper (Theorem~\ref{thm:polygon-einstein},
proved in Section~\ref{sec:lichnerowicz}) is:
under conditions (H1)--(H3), the polygon entropy functional
has a unique maximum at uniform curvature $R = \mathrm{const}$,
implying the vacuum Einstein equations $R_{\mu\nu} = \Lambda g_{\mu\nu}$.

\paragraph{The central charge.}
The quantity $c = 12\,b(N)$ is extracted from the Havelock
eigenvalue data as the unique mode-independent offset
(Proposition~\ref{prop:equiv-rr}).
It serves simultaneously as the WDW kinetic coefficient
(\S\ref{sec:wdw}), the CS level $k = c/6$
(\S\ref{sec:graviton}), the Brown--Henneaux central charge
($G := 3\ell/(2c)$), and the orbifold CFT ground-state
energy (\S\ref{sec:three-layer}).
If any two roles required different values of~$c$, the
framework would be falsified.

%% ====================================================================
\section{The graviton at $N = 7$ and the Chern--Simons lift}
\label{sec:graviton}

The stability boundary $N_{\mathrm{crit}} = 7$ has a
representation-theoretic explanation that connects
vortex dynamics to gravity through a chain of
proved identifications.

The linearised spectrum of the Thomson vortex system
coincides with that of the
Calogero--Moser--Sutherland (CMS) model at coupling $g = 1$
(Paper~I, Remark~3.1; \citealt{Calogero1975}).
The CMS model on a circle has dynamical symmetry algebra
$\mathfrak{sl}(2,\mathbb{R})$
\citep{OlshanetskyPerelomov1976}.
We now prove that the quadratic Casimir $C_2$ on the
$m$-th mode is $m(N{-}m)/2$, and that this is
the same $\mathfrak{sl}(2,\mathbb{R})$ that
appears in the CS formulation of gravity.

\begin{proposition}[CMS--CS Casimir identification]
\label{prop:casimir-havelock}
On $\mathbf{H}^2 = \mathrm{SL}(2,\mathbb{R})/\mathrm{SO}(2)$,
the $\mathfrak{sl}(2,\mathbb{R})$ Casimir restricted to the
$m$-th $\mathbb{Z}_N$ tangential mode of the regular $N$-gon is
\begin{equation}\label{eq:cms-casimir}
  C_2(m) = m(N{-}m)/2 = f(m,N).
\end{equation}
This holds simultaneously for the CMS dynamical symmetry
(Olshanetsky--Perelomov), the KZ connection (Etingof--Frenkel--Kirillov),
and the CS Wilson line (Witten), because all three use the
same $\mathfrak{sl}(2,\mathbb{R})$: the isometry algebra
of~$\mathbf{H}^2$.
\end{proposition}

\begin{proof}
Three steps: Laplacian, Havelock, identification.

\emph{Step 1 (Symmetric space theorem).}
On a Riemannian symmetric space $G/K$, the quadratic Casimir
of~$\mathfrak{g}$ acts as the Laplace--Beltrami operator
(\citealt{Helgason1984}, Chapter~II, Theorem~2.3).
For $G = \mathrm{SL}(2,\mathbb{R})$, $K = \mathrm{SO}(2)$,
$G/K = \mathbf{H}^2$:
the $\mathfrak{sl}(2,\mathbb{R})$ Casimir
$= -\Delta_{\mathbf{H}^2}$ (up to sign convention).

\emph{Step 2 (Restriction to the $N$-gon ring).}
The Laplace--Beltrami operator on $\mathbf{H}^2$ in
polar coordinates $(\rho, \theta)$ has angular part
$\Delta_\theta = (1/\sinh^2\!\rho)\,\partial^2/\partial\theta^2$.
At the $N$-gon ring of radius~$\rho_0$, the tangential Hessian
of $-\sum_{p<q}\log d_{\mathrm{hyp}}(z_p, z_q)$ is a
$\mathbb{Z}_N$-circulant with entries
$h_p = (4\sinh^2\!\rho_0\,\sin^2(\pi p/N))^{-1}$
(verified numerically to $10^{-15}$ for $N = 3,\ldots,12$
at five values of~$\rho_0$;
see Code Availability).

\emph{Step 3 (Havelock eigenvalue).}
The $\mathbb{Z}_N$ Fourier eigenvalue of this circulant is
\[
  h_m = \frac{1}{\sinh^2\!\rho_0}
  \sum_{p=1}^{N-1}\frac{1 - \cos(2\pi pm/N)}{4\sin^2(\pi p/N)}
  = \frac{m(N{-}m)/2}{\sinh^2\!\rho_0} = \frac{T_m}{\sinh^2\!\rho_0}
\]
by the Havelock identity (Paper~I, Lemma~\ref*{I-lem:havelock}).
In the conformal normalisation
(multiply by $\sinh^2\!\rho_0$, equivalently
use the Poincar\'e-disk coordinate $|z| = \tanh(\rho_0/2)$),
the Casimir eigenvalue is
$C_2(m) = T_m = m(N{-}m)/2$,
\emph{independent of the ring radius~$\rho_0$}:
it is a representation-theoretic invariant, not a geometric
quantity.
\end{proof}

\begin{remark}[Why CMS and CS have the same $\mathfrak{sl}(2,\mathbb{R})$]
\label{rmk:same-sl2}
The Olshanetsky--Perelomov $\mathfrak{sl}(2,\mathbb{R})$ (CMS
dynamical symmetry) and the Witten $\mathfrak{sl}(2,\mathbb{R})$
(CS gauge algebra) are a priori different algebras acting on
different spaces.
They coincide because both act as \emph{isometries of
the same $\mathbf{H}^2$}: the CMS model lives on the
configuration space
$\mathrm{SL}(2,\mathbb{R})/\mathrm{SO}(2)$
(the symmetric space whose Laplacian is the CMS Hamiltonian),
and the CS gauge group acts on the same
$\mathrm{SL}(2,\mathbb{R})/\mathrm{SO}(2)$ base manifold.
The Etingof--Frenkel--Kirillov theorem
(\citeyear{EtingofFrenkelKirillov1998}, Theorem~7.1.2)
makes this precise: the CMS Hamiltonian restricted to the
$\mathbb{Z}_N$-symmetric locus of the configuration space
\emph{is} the KZ connection (the flatness condition of
the CS gauge field with $N$ insertions).

The chain is:
\begin{enumerate}
\item \emph{CMS $\to$ KZ.}
  At the $\mathbb{Z}_N$-symmetric point, the CMS Hamiltonian
  $=$ the radial part of the KZ connection (EFK, Theorem~7.1.2).
\item \emph{KZ $\to$ CS Wilson line.}
  CS Wilson line expectation values
  $\langle W_m\rangle = \mathrm{Tr}_{R_m}
  \mathcal{P}\exp\!\oint\! A$
  satisfy the KZ equation \citep{Witten1988}.
\item \emph{No mode mixing.}
  At the $\mathbb{Z}_N$-fixed point, the KZ connection
  commutes with the $\mathbb{Z}_N$ action
  ($\mathbb{Z}_N$-equivariance at a fixed configuration).
  Each Fourier mode~$V_m$ is one-dimensional
  (for $m \ne N{-}m$): the $\mathbb{Z}_N$ irreducible
  with character $\omega^m$ ($\omega = e^{2\pi i/N}$)
  is $1$-dimensional over~$\mathbb{C}$; the palindromic
  pair $\{m, N{-}m\}$ shares a $2$-dimensional real
  eigenspace with the same eigenvalue.
  On a $1$D space, eigenvalue matching $\Leftrightarrow$
  operator matching.
\end{enumerate}
Therefore $C_2^{\mathrm{CMS}}(m)
= C_2^{\mathrm{KZ}}(m)
= C_2^{\mathrm{CS}}(m) = m(N{-}m)/2$.
\end{remark}

In the CS formulation of $2{+}1$D gravity
\citep{Witten1988}, the polygon is a Wilson line
whose eigenvalue is the $\mathrm{SL}(2,\mathbb{R})$ Casimir
$j(j{+}1)$ of the representation~$R_m$.
Setting $j(j{+}1) = m(N{-}m)/2$ determines $j$ for each mode.
The critical mode has $f(m^*,N) = j(j{+}1)$
with integer~$j$ only at specific~$N$:
\begin{center}
\small
\begin{tabular}{lccl}
\toprule
$N$ & $m^*$ & $f(m^*,N)$ & $j$ such that $j(j{+}1) = f$ \\
\midrule
$4$ & $2$ & $2$ & $j = 1$ (vector boson) \\
$7$ & $3$ & $6$ & $j = 2$ (graviton) \\
\bottomrule
\end{tabular}
\end{center}
\noindent
No other $N \le 30$ gives spin $j = 2$ at the critical mode.
The \emph{graviton} (spin-$2$ field of general relativity)
appears precisely at $N = 7$: the CMS/CS Casimir identification
places the $j = 2$ representation at the Havelock stability
threshold.
The total $\csc^2$ sum $\lambda_0 = \sum_{p=1}^{N-1}\csc^2(\pi p/N)
= (N^2{-}1)/3$ itself has the Casimir form
$\lambda_0 = (4/3)\,j(j{+}1)$ with $j = (N{-}1)/2$,
where $4/3 = 1 + 2B_2$ (Paper~I, \S\ref*{I-sec:b2-tower}).
The same $4/3$ that governs the self-interaction of the
$\csc^2$ kernel also sets the normalisation of the
total angular stiffness of the polygon.

\begin{proposition}[Graviton identification]
\label{prop:graviton}
$N = 7$ is the unique polygon number $N \ge 3$ for which the
critical Havelock Casimir $f(\lfloor N/2\rfloor, N)$
equals $j(j{+}1)$ with $j = 2$ (the graviton spin).
Higher Pell solutions ($N = 41$, $j = 14$; $N = 239$, $j = 84$; \ldots)
give $j > 2$ and lie outside the stability range $N \le 7$.
\end{proposition}

\begin{proof}
For even $N = 2k$: $f(k, 2k) = k^2/2$.
Then $j(j{+}1) = k^2/2$ gives
$j = (-1 + \sqrt{1 + 2k^2})/2$.
This is an integer iff $1 + 2k^2$ is a perfect square,
i.e.\ $n^2 - 2k^2 = 1$ (Pell equation).
Solutions: $(n,k) = (1,0), (3,2), (17,12), \ldots$
No even~$N$ yields $j = 2$: at $k = 2$ ($N = 4$), $j = 1$;
at $k = 12$ ($N = 24$), $j = 8$.
Setting $j = 2$ requires $n = 5$, giving $25 - 2k^2 = 1$,
i.e.\ $k^2 = 12$, which has no integer solution.

For odd $N = 2k{+}1$: $f(k, 2k{+}1) = k(k{+}1)/2$.
Then $j(j{+}1) = k(k{+}1)/2$ gives $j = (-1 + \sqrt{1+2k(k{+}1)})/2$.
At $k = 3$ ($N = 7$): $j(j{+}1) = 6$, $j = 2$.
At $k = 4$ ($N = 9$): $j(j{+}1) = 10$, $j \approx 2.70$ (not integer).
At $k = 5$ ($N = 11$): $j(j{+}1) = 15$, $j \approx 3.41$.
For $4 \le k \le 19$: $j$ is irrational (the next Pell
solution is $k = 20$, $N = 41$, $j = 14$).
Setting $y = 2j{+}1$ and $N = 2k{+}1$, the integer-spin
condition becomes the negative Pell equation $N^2 - 2y^2 = -1$.
The solutions are generated by $\varepsilon = 1 + \sqrt{2}$:
$(N,y) = (1,1),\, (7,5),\, (41,29),\, \ldots$,
giving $j = 0,\, 2,\, 14,\, \ldots$
The first non-trivial solution is $N = 7$, $j = 2$.
\end{proof}

\begin{remark}[Physical content of the operator coincidence]
\label{rmk:cms-cs-physical}
The numerical verification
($N = 3,\ldots,12$; see Code Availability)
confirms the eigenvalue identity $C_2^{\mathrm{CMS}} =
C_2^{\mathrm{KZ}} = C_2^{\mathrm{CS}} = m(N{-}m)/2$.
Operator coincidence on a mathematical space does not
by itself establish physical identification.
In $2{+}1$D, however, a stronger statement holds:
the Weyl tensor vanishes identically, so there are
\emph{no propagating gravitational modes}.
The only gravitational degrees of freedom are
\emph{boundary modes} (Brown--Henneaux 1986).
By the CS/WZW theorem \citep{Witten1988}, the boundary
modes of the CS gauge field are the WZW currents---which
are the polygon modes (Paper~IV, Theorem~\ref*{IV-thm:su3}, Step~6).
In $2{+}1$D, the boundary theory of CS gravity and the
Kirchhoff/Havelock vortex system are mathematically
identical: the CMS/CS operator coincidence is a
spectral isomorphism (eigenvalue matching on each mode).
\end{remark}

\paragraph{The Chern--Simons lift.}
In $2{+}1$ dimensions, gravity with $\Lambda < 0$ is a
Chern--Simons gauge theory with group
$\mathrm{SL}(2,\mathbb{R}) \times \mathrm{SL}(2,\mathbb{R})$
and level
\begin{equation}\label{eq:cs-level}
  k = \frac{\ell}{4G} = \frac{c}{6},
\end{equation}
where the Brown--Henneaux relation $c = 6k$ gives
$k = 2\,b(N)$ for the exact central charge $c = 12\,b(N)$
(with $b(N) = N(N{+}1)/12 - \ln 2 + (\ln N)/(N{-}1)$;
the proof via mean Casimir and the Gauss product
$\prod_{p=1}^{N-1}\sin(\pi p/N) = N/2^{N-1}$
is in Paper~I, \S5).
At leading order $c \approx N^2$, so $k \approx N^2/6$.
The CS level $k = c/6$ is not integer for most~$N$.
This is consistent for three independent reasons:
\begin{enumerate}
\item \emph{No topological obstruction.}
  Level quantisation for compact groups arises from
  $\pi_3(G) = \mathbb{Z}$ (the winding number of large gauge
  transformations on~$S^3$).
  For $\mathrm{SL}(2,\mathbb{R})$:
  $\pi_3(\mathrm{SL}(2,\mathbb{R})) = 0$
  \citep[\S3]{Witten1988}, so no integrality constraint exists.
\item \emph{Well-defined partition function.}
  The CS partition function on the Seifert manifold decomposes as
  $Z = e^{-kS_0} \times Z_{\mathrm{1\text{-}loop}} \times Z_{\mathrm{Wilson}}$,
  where only the classical factor $e^{-kS_0}$ depends on~$k$
  (the one-loop determinant and Wilson lines depend on the
  Havelock eigenvalues, which are geometric, not on~$k$).
  The classical factor is \emph{entire} in~$k$.
  The nearest singularity of $Z(k)$ (from the analytic
  continuation of the compact form, \citealt{Witten1991};
  \citealt{BarNatan1991}) is at $k = -h^\vee = -2$.
  Since $k = 2b(N) > 0$ for all $N \ge 3$, we are in the
  regular region
  (verified numerically for $N = 3,\ldots,15$;
  see Code Availability).
\item \emph{Semiclassical validity.}
  For $N \ge 7$: $k = 2b(N) > 8.6$, and
  $1/k$ corrections are $< 2\%$.
  For $N = 4$: $k = 2.3$, and $1/k$ corrections are
  ${\sim}\,25\%$---the semiclassical expansion is
  \emph{marginal} at the electroweak polygon.
  The exact (non-perturbative) CS partition function
  is used for quantitative predictions at $N = 4$.
\end{enumerate}
The compact gauge sectors ($\mathrm{SU}(2)_1$,
$\mathrm{SU}(3)_1$) have integer levels;
only the gravitational $\mathrm{SL}(2,\mathbb{R})$ sector
carries the irrational level $k = 2b(N)$, which is
well-defined for the non-compact group.
The vortex polygon is a \emph{Wilson line}
$W_m = \mathrm{Tr}_{R_m}\,\mathcal{P}\exp\!\oint\! A$
in the representation $R_m$ with Casimir $f(m,N)$.
The Havelock eigenvalue is the Wilson line expectation
$\langle W_m \rangle = e^{-\lambda_m L}$; the three-layer
decomposition corresponds to
bulk (cosmological) $\times$ representation (matter) $\times$
one-loop (quantum) factors.
The CS partition function on Seifert manifolds is computed
via Stieltjes--Wigert polynomials \citep{Tierz2006},
linking the Havelock eigenvalue sequence to the
log-normal moment problem.

(The Havelock matrix is a circulant on $\mathbb{Z}/N\mathbb{Z}$
with the log-sine kernel; the Casimir $m(N{-}m)/2$ is the
Fourier eigenvalue and $\delta_m$ the lattice correction.)

\section{The spectral bound as a quantum gravitational bound}
\label{sec:spectral-bound}

The spectral correction and the gravitational one-loop
partition function are both moments of the heat kernel:
\begin{alignat}{2}
  \log Z_{\mathrm{grav}}
    &= -\int_0^\infty
      \bigl[\operatorname{Tr}\,e^{-t\Delta_S} - 1\bigr]\,
      \frac{dt}{t}
    &&\qquad\text{(weight $1/t$),}
    \label{eq:grav-moment} \\
  \delta C_1(z_0)
    &= \int_0^\infty
      \bigl[K(z_0,z_0,t) - K_{\mathbf{H}^2}(t)\bigr]\,dt
    &&\qquad\text{(weight $1$).}
    \label{eq:vortex-moment}
\end{alignat}
The first is the functional determinant $\det\Delta_S$,
the second determines polygon stability.
Both involve the Selberg zeta function $Z_\Gamma(s)$
at the zero $s = 1$, but this requires care.

\paragraph{The zero at $s = 1$.}
The Selberg zeta $Z_\Gamma(s)$ has zeros at $s$ such that
$s(1{-}s)$ is an eigenvalue of $\Delta_S$.
The constant eigenfunction ($\lambda_0 = 0$) gives a zero
at $s = 1$, so $Z_\Gamma(1) = 0$---not a finite number.
The Euler product
$\prod_\gamma \prod_k (1 - e^{-(s+k)\ell_\gamma})$
converges for $\mathrm{Re}(s) > 1$ but diverges at $s = 1$
(by the prime geodesic theorem: the partial product over
geodesics with $\ell \le L$ decays to~$0$ as $L \to \infty$).

The physically relevant quantities are the \emph{residue}
at this zero:
\begin{itemize}
  \item $Z_\Gamma'(1)$: the \emph{regularised} determinant
    $\det{}'\Delta_S = \prod_{n \ge 1} \lambda_n$ (omitting
    $\lambda_0 = 0$).
    This is the gravitational one-loop partition function
    \citep{GiombiMaloneyYin2008}.
  \item The spectral sum
    $\sum_{n \ge 1} |\varphi_n(z_0)|^2 / \lambda_n$:
    the regularised Green's function at the diagonal,
    which determines $\delta C_1(z_0)$.
\end{itemize}
Both are finite and well-defined; both involve the
analytic structure of $Z_\Gamma$ at its zero, not the
(divergent) Euler product.

\paragraph{The spectral computation of $\delta C_1$.}
At the Aut-fixed centre~$z_0$ of the Bolza surface,
only trivial-representation eigenfunctions contribute:
$\delta C_1(z_0) = \sum_{n:\,\text{triv}}
|\varphi_n(z_0)|^2 / \lambda_n \approx 0.011$,
dominated by the first trivial-rep eigenvalue
$\lambda_{\mathrm{triv}} \approx 15.05$
(the first eigenvalue in the trivial representation of
$\mathrm{Aut}(\mathrm{Bolza})$; the full spectral gap is
$\lambda_1 \approx 3.84$;
\citealt{Strohmaier2013}), with rapid convergence:
$1/(4\pi \cdot 15.05) \approx 0.005$,
$1/(4\pi \cdot 23.08) \approx 0.003$,
$1/(4\pi \cdot 32.67) \approx 0.002$, $\ldots$.
The first~$8$ trivial-representation eigenvalues of
\citet{Strohmaier2013} give
$\sum_{n=1}^{8} 1/(4\pi\lambda_n) = 0.019$.
The tail is bounded by Sogge's sup-norm estimate
($|\phi_n(z_0)|^2 \le C\lambda_n^{1/2}$)
combined with Weyl asymptotics ($\lambda_n \ge c\,n$
for the trivial representation):
$\sum_{n > 8} 1/(4\pi\lambda_n^{3/2}) < 0.001$.
Therefore
$|\delta C_1| \le 0.019 + 0.001 = 0.020$.

The bound $|\delta C_1| \le 0.020$ on the Bolza surface
is conditional on the Strohmaier--Uski spectral data
(numerical computation, not a closed-form proof)
and unconditional with respect to the Selberg
eigenvalue conjecture (the bound uses only
$\lambda_{\mathrm{triv}} \approx 15.05$,
well above the Selberg threshold $1/4$).
The spectral gap $\lambda_1$ controls both $\delta C_1$
(through the sum $\sum 1/\lambda_n$) and the gravitational
partition function (through $\det{}'\Delta_S$), giving a
\emph{structural parallel} between vortex stability and
quantum gravity---but not, at present, a sharp functional
inequality linking the two.
In the gravitational language:
\begin{itemize}
  \item $C_1(\mathbf{H}^2, \xi)$: classical stability
    (flat-space AdS, no topology).
  \item $\delta C_1$: one-loop quantum correction
    (from the topology of the spatial slice).
  \item The bound $|\delta C_1| \le 2.86/(4\pi\lambda_{\mathrm{triv}})$
    at the Aut-fixed centre
    ($\lambda_{\mathrm{triv}} \approx 15.05$ for the Bolza surface),
    or $|\delta C_1| \le 2.86/(4\pi\lambda_1)$ at a generic point
    ($\lambda_1 \approx 3.84$): quantum gravitational corrections
    to polygon stability.
\end{itemize}

\paragraph{The Selberg dictionary (structural parallel, not a sharp inequality).}
Four distinct physical quantities involve the Selberg
zeta at the zero $s = 1$.
The entries below are structural analogies; no functional inequality
linking the two sides is claimed.
\begin{center}
\small
\begin{tabular}{ll}
\toprule
$Z_\Gamma$ data at $s = 1$ & Physical meaning \\
\midrule
$Z_\Gamma'(1)$ &
  Gravitational partition function \\
  \quad(regularised determinant $\det{}'\Delta_S$) &
  $= \prod_{n \ge 1}\lambda_n$ \\[4pt]
$\sum_{n \ge 1} |\varphi_n(z_0)|^2/\lambda_n$ &
  Vortex stability correction $\delta C_1$ \\
  \quad(spectral; convergent) &
  Localised at $z_0$; Aut-suppressed \\[4pt]
Palindromic threshold polynomials &
  Stability transitions: values of $\xi$ \\
  \quad(from $C_1(\xi) = f_{\max}(N)$) &
  where $N_{\mathrm{crit}}$ jumps \\[4pt]
Hecke eigenvalues of associated forms &
  Prime-by-prime visibility \\
  \quad($a_p$ from modular forms) &
  Which thresholds are visible on $\Gamma(p)$ \\
\bottomrule
\end{tabular}
\end{center}
The first two rows involve the same zero $s = 1$ of
$Z_\Gamma$ but differ in what they extract:
the determinant (a \emph{product} of eigenvalues) versus
the resolvent trace (a \emph{sum} of $1/\lambda_n$,
localised at~$z_0$).
Neither can be computed from the Euler product
(which diverges at $s = 1$); both require either the
spectral data or the analytic continuation of $Z_\Gamma$.

The palindromic thresholds control
\emph{where} stability transitions occur (at which
curvature $\xi$); the eigenvalue spectrum controls
\emph{how large} the correction $\delta C_1$ is.
On highly symmetric surfaces (the modular surface,
the Bolza surface), the systole and the palindromic
geodesics coincide; on a generic arithmetic surface
they are distinct.

\section{Holographic interpretation}
\label{sec:holographic}

In AdS$_3$/CFT$_2$, a point mass $m$ in the bulk creates a
conical deficit that corresponds to the insertion of a heavy
operator $\mathcal{O}$ of conformal dimension
$h = (c/24)(1 - (1 - 8Gm)^2)$ in the boundary CFT
\citep{Carlip1998}.
A regular $N$-gon of equal masses at radial coordinate~$r$
in the bulk corresponds to $N$ equally-spaced heavy operator
insertions on the boundary circle.

In this language, the Havelock decomposition takes the form:
\begin{center}
\begin{tabular}{ll}
\toprule
Bulk (2+1D AdS gravity) & Boundary (1+1D CFT) \\
\midrule
Perturbation mode $m$ &
  $\mathbb{Z}_N$ conformal block channel \\
$f(m,N) = m(N{-}m)/2$ &
  Casimir of the $m$-th channel \\
$C_1(\xi)$ &
  Function of $h/c$ and AdS radius \\
$\lambda_m = C_1 - f(m,N) > 0$ &
  Positivity of the $m$-th block \\
$C_1 \ge f_{\max}(N)$ &
  Bootstrap-type bound on $h/c$ \\
\bottomrule
\end{tabular}
\end{center}
The stability condition $C_1 \ge m(N{-}m)/2$ for all~$m$
becomes a \emph{lower bound on the conformal dimension} of the
boundary operators: below this bound, the bulk polygon is
unstable and the corresponding $\mathbb{Z}_N$ conformal block
changes sign.
The palindromic threshold $\xi^*(N)$ translates to a critical
ratio $h/c$ at which this sign change occurs.

The kinematic part of this dictionary is confirmed by the
following identity.

\begin{proposition}[Havelock Casimir = one-chiral OPE Casimir]
\label{prop:chiral-casimir}
Let $T_m = \sum_{p=1}^{N-1}
  \frac{1 - \cos(2\pi pm/N)}{2\sin^2(\pi p/N)}$
be the $\mathbb{Z}_N$ Fourier eigenvalue of the OPE
kernel $\csc^2(\pi p/N)$.  Then
\begin{equation}\label{eq:chiral-ope}
  T_m = m(N{-}m) = 2\,f(m,N),
\end{equation}
exactly twice the Havelock Casimir $f(m,N) = m(N{-}m)/2$.
\end{proposition}

\begin{proof}
Write $\omega = e^{2\pi i/N}$ and
$S_m = \sum_{p=1}^{N-1}(1 - \omega^{pm})/|1 - \omega^p|^2$.
Since $|1 - \omega^p|^2 = 2(1 - \cos(2\pi p/N))
= 4\sin^2(\pi p/N)$ and
$1 - \omega^{pm} = 1 - \cos(2\pi pm/N) - i\sin(2\pi pm/N)$,
the real part gives
$T_m = \mathrm{Re}(S_m)
= \sum_{p=1}^{N-1}[1 - \cos(2\pi pm/N)]/(2\sin^2(\pi p/N))$.
Now consider the polynomial identity on
$\mathbb{Z}/N\mathbb{Z}$: the function
$f_m(p) = [1 - \cos(2\pi pm/N)]/(2\sin^2(\pi p/N))$
extends to a trigonometric polynomial of degree~$\le N{-}1$.
Its sum equals the coefficient of $z^0$ in the Laurent
expansion of $\frac{1 - z^m}{1 - z}\cdot\frac{1 - z^{-m}}{1 - z^{-1}}$,
evaluated at $z = \omega$ and summed over $p = 1,\ldots,N{-}1$.
This product is $|\sum_{k=0}^{m-1}z^k|^2$, whose
$z^0$ coefficient (from the orthogonality
$\sum_{p=0}^{N-1}\omega^{pn} = N\delta_{n,0}$) is
$\min(m, N{-}m) \cdot \max(m, N{-}m) = m(N{-}m)$,
since exactly $m$ terms contribute in each sector.
\end{proof}

The factor of~$2$ has physical content.
The logarithmic interaction decomposes as
$\log|z|^2 = \log z + \log \bar{z}$
(holomorphic $+$ anti-holomorphic).
In the boundary CFT, the OPE of two heavy operators has
both chiral sectors, giving the \emph{full} Casimir $m(N{-}m)$.
The Havelock eigenvalue uses $f = m(N{-}m)/2$, the
contribution of \emph{one} chiral sector, because:
\begin{itemize}
  \item The Kirchhoff equations $\dot{z}_k = -i\,\partial H/\partial\bar{z}_k$
    are first-order: coordinates are conjugate to each other
    (no separate momenta).
    This halves the phase-space dimension, selecting one chiral
    sector of the $\log|z|^2 = \log z + \log\bar{z}$ interaction.
  \item The stability criterion $\lambda_m > 0$ is
    $C_1 > m(N{-}m)/2$: the confining potential must exceed
    the holomorphic OPE repulsion in the $m$-th
    $\mathbb{Z}_N$ channel.
\end{itemize}

In gravity (second-order dynamics), both sectors contribute
to the equations of motion, but the stability criterion---
positivity of the Hessian eigenvalue---still involves
$m(N{-}m)/2$ per sector
(Proposition~\ref{II-prop:grav-stability}).

The remaining step is to express $C_1$ in terms of boundary
CFT data.
On the flat plane, $C_1 = N{-}1$ (the Havelock self-energy);
on $\mathbf{H}^2$ at curvature parameter~$\xi$,
\begin{equation}\label{eq:C1-H2}
  C_1(\xi) = (N{-}1)\,\frac{1{+}\xi^2}{(1{-}\xi)^2}.
\end{equation}
The curvature parameter $\xi$ encodes the radial
position of the bulk polygon; through the Brown--Henneaux
relation $c = 3\ell/(2G)$ and the deficit angle--dimension
map $h = (c/24)(1{-}(1{-}8Gm)^2)$, the stability condition
becomes a bound on conformal dimensions parametrised by the
bulk position.
The palindromic thresholds $\xi^*(N)$ give algebraic-number
bulk positions at which the stability bound is saturated:
\begin{center}
\small
\begin{tabular}{cccccl}
\toprule
$N$ & $\xi^*(N)$ & $d_{\mathrm{crit}}/\ell$ &
  $\varepsilon = 1{-}\xi^*$ & $h/c_{\max}$ & Field \\
\midrule
$8$  & $8{-}3\sqrt{7}$            & $0.51$ & $0.937$ & $0.018$
  & $\mathbb{Q}(\sqrt{7})$ \\
$10$ & $1/7$                       & $0.80$ & $0.857$ & $0.015$
  & $\mathbb{Q}$ \\
$12$ & $(18{-}5\sqrt{11})/7$       & $0.97$ & $0.798$ & $0.013$
  & $\mathbb{Q}(\sqrt{11})$ \\ [2pt]
$23$ & $\varphi^{-2}$              & $1.44$ & $\varphi^{-1}$ & $0.007$
  & $\mathbb{Q}(\sqrt{5})$ \\
\bottomrule
\end{tabular}
\end{center}
Here $d_{\mathrm{crit}}/\ell$ is the geodesic distance from
the AdS centre in units of the AdS radius, and
$\varepsilon = 1 - \xi^*$ is the boundary proximity
(the polygon must sit at $\varepsilon < \varepsilon^*(N)$,
i.e., sufficiently close to the AdS boundary, to be stable).
The $N = 23$ threshold gives
$\varepsilon^* = \varphi^{-1} = (\sqrt{5}{-}1)/2$: the golden
ratio appears as a critical bulk-position parameter in the
holographic dictionary.

The two constraints on the polygon are independent:
\begin{itemize}
  \item \emph{Mass}: $h/c < h/c_{\max}(N)$
    (total deficit $< 2\pi$; from the bulk geometry).
  \item \emph{Position}: $\xi > \xi^*(N)$
    (confining potential $> $ OPE repulsion; from the
    Havelock decomposition and the chiral
    identity~\eqref{eq:chiral-ope}).
\end{itemize}
Together, they define the allowed region in the
$(h/c,\;\xi)$ parameter space for a stable $N$-gon of
heavy operators---a two-parameter family of bootstrap-type
bounds, one for each $N$, all descending from the single
identity $\lambda_m = C_1 - m(N{-}m)/2$.

\paragraph{The confining potential from the OPE kernel.}
The formula $C_1 = (N{-}1)(1{+}\xi^2)/(1{-}\xi)^2$ can be
derived either from the Kimura Hamiltonian for point vortices
on $\mathbf{H}^2$ (Paper~I, \S4) or from the second radial
variation of the chordal OPE kernel
$\sigma(z,w) = |z{-}w|/|1{-}\bar{z}w|$ at the $N$-gon.
The two derivations are mathematically identical:
both compute the Hessian of $-\sum_{j<k}\log\sigma(z_j,z_k)$
on the Poincar\'e disk.
The factor $(1{+}\xi^2)/(1{-}\xi)^2$ is \emph{not} the
conformal factor $\Omega(\sqrt{\xi})^2 = 1/(1{-}\xi)^2$ alone;
the additional $(1{+}\xi^2)$ comes from the angular structure
of $\sigma$ under radial perturbations.

\paragraph{Relation to the conformal bootstrap.}
The stability condition $\lambda_m > 0$ and the conformal
bootstrap positivity conditions (crossing symmetry $+$
unitarity, \citealt{Rattazzi2008}) are both positivity
conditions on conformal block coefficients, but in
\emph{different channels}:
\begin{center}
\begin{tabular}{lll}
\toprule
& Bootstrap & Vortex stability \\
\midrule
What is constrained &
  CFT spectrum $\{\Delta_i\}$ &
  Bulk polygon position $\xi$ \\
Input &
  $c$, $\Delta_{\mathrm{gap}}$ &
  $N$, $h/c$ \\
Positivity condition &
  OPE $f_{ijk}^2 \ge 0$ &
  $\lambda_m = C_1 - f(m,N) > 0$ \\
Channel &
  generic 4-point OPE &
  $\mathbb{Z}_N$-symmetric \\
\bottomrule
\end{tabular}
\end{center}
The two are \emph{complementary}: the bootstrap constrains
which CFTs exist, while the vortex stability constrains
which bulk configurations are stable in a given CFT.
The identification $\lambda_m \leftrightarrow$
OPE coefficient would make the stability condition a
bootstrap condition for $\mathbb{Z}_N$-symmetric
correlators---a direction requiring the $\mathbb{Z}_N$
conformal block computation sketched below.

\section{The $\lambda_m \leftrightarrow$ OPE identification}
\label{sec:ope-identification}

The Havelock eigenvalue is entirely computable from boundary
CFT data.
The key is to decompose the semiclassical $N$-point function
$\langle\mathcal{O}(z_1)\cdots\mathcal{O}(z_N)\rangle$
of $N$ equally-spaced heavy operators into
$\mathbb{Z}_N$ channels: the angular channels give $m(N{-}m)/2$
(the kinematic OPE Casimir, Proposition~\ref{prop:chiral-casimir}),
and the radial confinement gives $C_1(\xi)$.

In the semiclassical (large-$c$) limit, the $N$-point
function factorises as
\begin{equation}\label{eq:n-point-semiclassical}
  \log\langle\mathcal{O}(z_1)\cdots\mathcal{O}(z_N)\rangle
  \;\approx\;
  -2h \sum_{j<k} \log\sigma(z_j, z_k)
  \;+\; O(1),
\end{equation}
where $\sigma = |z_j{-}z_k|/|1{-}\bar{z}_j z_k|$ is the
chordal distance.
The leading term is the ``free OPE energy''; the $O(1)$
correction comes from the Virasoro block (graviton exchange
in the bulk).

Under \emph{angular} perturbation
$z_k \to R\,e^{i(2\pi k/N + \delta\theta_k)}$ at the $N$-gon,
the second variation of the leading term decomposes into
$\mathbb{Z}_N$ Fourier modes $\delta\theta_k =
\sum_m a_m \cos(2\pi mk/N)$:
\begin{equation}\label{eq:second-var-ope}
  \delta^2\!\Bigl[-2h\sum_{j<k}\log\sigma\Bigr]
  = 2h \sum_{m=1}^{N-1} m(N{-}m)\,|a_m|^2.
\end{equation}

\begin{proposition}[Angular OPE stiffness]
\label{prop:ope-stiffness}
The coefficient of $|a_m|^2$
in~\eqref{eq:second-var-ope} is $2h\cdot m(N{-}m)
= 2h \cdot T_m$, where $T_m = m(N{-}m)$ is the full
OPE Casimir (Proposition~\ref{prop:chiral-casimir}).
This follows analytically from
Proposition~\ref{prop:chiral-casimir}
($T_m = m(N{-}m)$, proved) and the constrained second
variation; verified numerically for $N \le 12$.
\end{proposition}

The missing piece is the confining potential $C_1(\xi)$.
We now show it too arises from the boundary.

\begin{theorem}[$\mathbb{Z}_N$ conformal block $=$ Havelock eigenvalue]
\label{thm:block-havelock}
Let $\mathcal{O}_1, \ldots, \mathcal{O}_N$ be $N$ identical heavy
operators of dimension $h$, placed at equally-spaced points
$z_k = \sqrt{\xi}\,e^{2\pi ik/N}$ on the boundary of the
Poincar\'e disk.
In the semiclassical limit ($c \to \infty$, $h/c$ fixed):
\begin{enumerate}
\item[\emph{(a)}] \emph{Angular stiffness.}
  The $m$-th $\mathbb{Z}_N$ channel of the angular second
  variation of the $N$-point function gives
  $2h\cdot m(N{-}m)$
  (Proposition~\ref{prop:ope-stiffness}).
\item[\emph{(b)}] \emph{Radial confinement.}
  The log-radial ($u = \log\!\sqrt{\xi}$) second variation
  of $-\sum_{j<k}\!\log\sigma(z_j, z_k)$ at the $N$-gon
  gives the mode-independent confining potential
  $C_1(\xi) = (N{-}1)(1{+}\xi^2)/(1{-}\xi)^2$.
\item[\emph{(c)}] \emph{Full identification.}
  The constrained second variation
  (angular $+$ radial at fixed angular impulse), divided
  by $2h$ per chiral sector, is
  \begin{equation}\label{eq:block-equals-havelock}
    \lambda_m \;=\;
    \underbrace{C_1(\xi)}_{\text{radial confinement}}
    \;-\;
    \underbrace{m(N{-}m)/2}_{\text{one-chiral OPE repulsion}}.
  \end{equation}
  This is the Havelock eigenvalue, derived entirely from
  boundary data.
\end{enumerate}
\end{theorem}

\begin{proof}
Part~(a) is Proposition~\ref{prop:ope-stiffness}.
For part~(b): the log-radial second variation of the
hyperbolic energy gives $C_1(\xi)$ as in~\eqref{eq:C1-H2}
(Paper~I, equation~\ref*{I-eq:C1_H2})
(the Riemannian Havelock identity on~$\mathbf{H}^2$).
Part~(c) combines (a) and (b) via the standard
Lagrange-multiplier argument (Paper~I, \S3).

\emph{Identity block dominance for $N > 4$.}
The theorem requires that the identity Virasoro block dominates the
OPE decomposition in each $\mathbb{Z}_N$ channel.
For $N = 4$ this is standard (the semiclassical $4$-point block is
the sole saddle on $\mathcal{M}_{0,4} \cong \mathbb{P}^1$).
For $N > 4$: the $\mathbb{Z}_N$-symmetric $N$-point correlator of
identical operators on $\mathbb{P}^1$ (\emph{genus zero})
equals the $1$-point function of the twist field~$\sigma_m$
on the orbifold $\mathbb{C}/\mathbb{Z}_N$
\citep{DixonHarveyVafaWitten1985}.
The orbifold has twisted-sector vacua $|k\rangle$ with
conformal dimensions $h_k = k(N{-}k)/(2N^2) \times c > 0$
for $k = 1, \ldots, N{-}1$.
On genus zero, these contribute to the $1$-point function
only as virtual intermediate states, suppressed by
$O(e^{-2\pi h_k}) = O(e^{-c})$ (exponentially small at
large~$c$, since $h_k \propto c$).
The identity block gives the leading (classical) contribution
$\exp(-c \times S_{\mathrm{conical}})$;
the $\mathbb{Z}_N$ symmetry reduces the moduli space
$\mathcal{M}_{0,N+1}^{\mathbb{Z}_N}$ to dimension~$2$
(ring radius and overall phase).
On this $2$D space, the identity block dominates
by two independent bounds:
\begin{enumerate}
\item \emph{Eigenvalue bound
(Proposition~\ref{prop:non-crossing}):}
  the $O(1/c)$ Virasoro corrections satisfy
  $|\delta f_m/\lambda_m| < 1/4$ for all non-critical~$m$,
  so no eigenvalue changes sign at finite~$c$
  and the symmetric saddle remains non-degenerate.
\item \emph{OPE coefficient bound (Liouville action):}
  the $\mathbb{Z}_N$ twist field $\sigma_k$ of order~$k$
  lives on the orbifold $\mathbb{C}/\mathbb{Z}_N$
  with conical metric of opening angle $2\pi/N$.
  Its OPE coefficient is
  $C_k^2 = \exp(-S_L(\mathrm{cone}))$
  where $S_L$ is the classical Liouville action
  \citep{DixonHarveyVafaWitten1985,Zamolodchikov1987}.
  The Liouville action on a $\mathbb{Z}_N$ cone regularised
  by the Gauss log-sine product (the same regularisation
  that gives $b(N)$ in the Havelock decomposition,
  Paper~I, \S\ref*{I-sec:algebraic-structure}) is
  $S_L \ge b(N)\ln N$.
  Therefore $C_k^2 \le \exp(-b(N)\ln N)$, and the
  total non-identity contribution is
  \[
    \frac{\delta F}{F}
    \le (N{-}1)\exp\!\bigl(-b(N)\ln N\bigr)
  \]
  (verified numerically for $N = 3,\ldots,16$;
  see Code Availability):
  $11\%$ at $N = 5$,
  $1.7\%$ at $N = 6$,
  $0.14\%$ at $N = 7$,
  $\sim 10^{-10}$ for $N = 11$.
\end{enumerate}
Bound~(1) excludes competing saddles at the eigenvalue
level; bound~(2) controls the non-perturbative twist
contributions to the conformal block decomposition.
Together they establish identity block dominance for
$N \ge 6$ (corrections $\le 2\%$).
At $N = 5$, the $11\%$ correction is the leading
systematic uncertainty; however, since
$\lambda_{\min}(5) = 1$ on the flat plane, a shift of
$\sim\!0.11$ leaves $\lambda_{\min} > 0$, so the
stability conclusion is not affected.
\end{proof}

Theorem~\ref{thm:block-havelock} identifies
$\lambda_m$ as the $m$-th $\mathbb{Z}_N$-channel stiffness
of the semiclassical $N$-point function.
No bulk computation is needed:
the stability condition $\lambda_m > 0$ is
\emph{the confining potential exceeding the holomorphic OPE
repulsion in channel~$m$}, a statement in the boundary CFT.
The palindromic thresholds
$\xi^*(N)$---at which $C_1(\xi^*) = f_{\max}(N)$---are
visible as algebraic-number bulk positions where the
$\mathbb{Z}_N$ conformal block changes sign.

\paragraph{Quantum correction.}
The $\mathbb{Z}_N$ Fourier decomposition of the Virasoro block
at central charge $c$ confirms the identification at all orders.
The $m$-dependent part of $\log F^{(m)}$ is
\begin{equation}\label{eq:block-1-over-c}
  \biggl[\frac{m(N{-}m)}{2}
         - \frac{m^2(N{-}m)^2}{4c}
         + O(1/c^2)\biggr]\,\log q.
\end{equation}
At $c = \infty$: the leading term is the Havelock Casimir (exact,
verified for all $N = 4,\ldots,12$ and all modes).
At finite~$c$: the correction $\delta h_m = -f_m^2/c$ is a
\emph{polynomial in the Casimir}, so the algebraic structure
(palindromic symmetry, number field) is preserved at each order
in~$1/c$.  The fractional correction $|f_m/c| \le 1/8$
(maximised at $m = N/2$, $c = N^2$) does not change any eigenvalue
sign (Proposition~\ref{prop:non-crossing}).


\section{The polygon--BTZ phase transition}
\label{sec:btz}

The polygon framework \emph{must} reproduce the known
Hawking--Page phase transition of AdS$_3$ gravity
\citep{Witten1998AdS}; this section verifies that it does.
The Euclidean action comparison between the polygon and the
BTZ black hole makes the Hawking--Page language precise
and gives a novel $N$-dependent correction
$T_c(N) < T_{\mathrm{HP}}$ from the polygon binding energy.

\paragraph{The Euclidean actions.}
At inverse temperature $\beta$, the regularised Euclidean
actions of the competing configurations are:
\begin{align}
  I_{\mathrm{polygon}}(\beta)
    &= -\frac{\beta}{8G} + \beta E_{\mathrm{polygon}},
    \label{eq:I-polygon} \\
  I_{\mathrm{BTZ}}(\beta)
    &= -\frac{\pi^2 \ell^2}{2G\beta}
    \qquad (M = 0\text{ threshold BTZ}),
    \label{eq:I-btz}
\end{align}
where $E_{\mathrm{polygon}} = Nm(1{-}2Gm) + V_{\mathrm{int}}$
is the total energy of the $N$-gon (kinematic mass $+$
gravitational binding), and $V_{\mathrm{int}} < 0$ is the
pairwise Havelock interaction on $\mathbf{H}^2$.
The vacuum action $-\beta/(8G)$ is the standard thermal-AdS$_3$
contribution.

\paragraph{The transition temperature.}
Setting $I_{\mathrm{polygon}} = I_{\mathrm{BTZ}}$:
\begin{equation}
  \label{eq:polygon-hp}
  T_c(N,m,\xi) \;=\;
    \frac{1}{2\pi\ell}
    \sqrt{1 - 8G\,E_{\mathrm{polygon}}(N,m,\xi)},
\end{equation}
which reduces to the standard Hawking--Page temperature
$T_{\mathrm{HP}} = 1/(2\pi\ell)$ for $E = 0$.
For $T < T_c$: the polygon has lower action
($\Delta I < 0$, polygon preferred).
For $T > T_c$: the BTZ has lower action.
The total polygon energy $E_{\mathrm{polygon}}$
lowers $T_c$ below $T_{\mathrm{HP}}$: more mass in the
polygon makes the BTZ transition harder, because the
binding energy deepens the polygon potential well.

\paragraph{The negative mode.}
For $N > N_{\mathrm{crit}}(\xi)$: the polygon has a negative
Havelock eigenvalue
$\lambda_m < 0$ at mode $m = \lfloor N/2 \rfloor$
(\S\ref{sec:lichnerowicz}).
By the vortex-gravity isomorphism
(Paper~II, Proposition~\ref{II-prop:grav-stability}),
the Havelock Hessian of the polygon is the second variation
of the gravitational energy around the multi-cone saddle.
A negative Havelock eigenvalue therefore is a negative mode
of the Lichnerowicz operator $\Delta_{\mathrm{Lich}}$ on the
multi-cone geometry---a perturbation that lowers the Euclidean
gravitational action, signaling decay
\citep{GiombiMaloneyYin2008}.
The mode $m = \lfloor N/2\rfloor$ displaces adjacent masses
alternately inward and outward, a pattern that drives pairwise
coalescence.
In the nonlinear regime, the coalescing masses form a single
body exceeding the BTZ threshold $8GM > 1$, completing the
decay to a BTZ black hole.

\paragraph{The three-region phase diagram.}
The polygon--BTZ competition has three distinct regimes:
\begin{enumerate}
\item \emph{Stable polygon, preferred}
  ($N \le N_{\mathrm{crit}}$, $T < T_c$):
  all Havelock eigenvalues are positive and the polygon has
  lower Euclidean action than the BTZ.
  The polygon contributes as a stable saddle to the path
  integral.
\item \emph{Stable polygon, metastable}
  ($N \le N_{\mathrm{crit}}$, $T > T_c$):
  all eigenvalues are positive but the BTZ has lower action.
  The polygon is a local minimum separated from the BTZ by a
  barrier; the tunnelling rate is $\Gamma \sim e^{-S_{\mathrm{BTZ}}}$.
\item \emph{Unstable polygon, decay}
  ($N > N_{\mathrm{crit}}$):
  the mode $m = \lfloor N/2\rfloor$ is a negative mode, and
  the polygon decays continuously (no barrier) toward the BTZ.
  The palindromic threshold $\xi^*(N)$ marks the boundary
  between regions~1/2 and region~3.
\end{enumerate}

\paragraph{The critical exponent.}
At the palindromic threshold $\rho = \rho^*(N)$, the critical mode
$m^* = \lfloor N/2\rfloor$ vanishes linearly:
$\lambda_{m^*}(\rho) \sim \coth(\rho^*) \cdot (\rho - \rho^*)$.
The one-loop partition function (Gaussian integration over the
zero mode) gives the universal critical exponent $\nu = 1/2$:
\[
  Z(\rho) \;\sim\; \frac{Z_{\mathrm{frozen}}}{|\rho - \rho^*|^{1/2}},
  \qquad
  Z_{\mathrm{frozen}} = \frac{2^{3(N-2)/4}}{(N{-}3)!!}
\]
where $(N{-}3)!! = 1 \cdot 3 \cdot 5 \cdots (N{-}3)$ is the double
factorial (even~$N$, one zero mode at $m^* = N/2$;
for odd~$N$, the palindromic pair $m^*$ and $N{-}m^*$ both
have $\lambda = 0$, giving two zero modes and
$Z \sim |\rho{-}\rho^*|^{-1}$.
$Z_{\mathrm{frozen}}$ for odd~$N$ is
$\prod_{m \ne m^*,\, m \ne N{-}m^*}
[f(m^*,N) - f(m,N)]^{-1/2}$,
computed from the Casimir gaps
$f(m^*,N) - f(m,N) = [m^*(N{-}m^*) - m(N{-}m)]/2$).
The closed-form expression includes the normalisation of the
first-order Kirchhoff path integral measure; it is verified
numerically against the gap-by-gap product for all~$N \le 20$
(see Code Availability).
The polygon--BTZ transition is second-order with mean-field exponents.

\paragraph{The entropy gap.}
In the microcanonical ensemble (fixed total energy),
the BTZ \emph{always} has higher entropy:
$\Delta S = S_{\mathrm{BTZ}}(E_{\mathrm{polygon}}) > 0$.
The polygon is therefore always entropically disfavored;
its survival as a metastable state depends entirely on the
positivity of the Havelock eigenvalues (the energy barrier),
not on entropy.
This makes $N_{\mathrm{crit}}$ the sharp dividing line:
below it, the polygon is a long-lived metastable state;
above it, the polygon decays on a dynamical timescale
$\tau \sim 1/\sqrt{|\lambda_{\lfloor N/2\rfloor}|}$.

\paragraph{Microscopic entropy accounting.}
The BTZ entropy decomposes into discrete and continuous contributions:
\begin{equation}\label{eq:btz-entropy-decomposition}
  S_{\mathrm{BTZ}}(N)
  = \underbrace{S_{\text{1-loop}}(N)}_{\text{continuous}}
  + \underbrace{\Delta S_{\text{wc}}}_{\text{discrete}},
\end{equation}
where $S_{\text{1-loop}} = \log(N{-}3)!! - \tfrac{3(N{-}2)}{4}\log 2$
(from the Casimir gaps) and $\Delta S_{\text{wc}}$ is the wall-crossing
contribution.
At each palindromic threshold, the $\mathbb{Z}/2$ Galois action
($m \leftrightarrow N{-}m$) contributes:
\begin{center}
\small
\begin{tabular}{lcl}
\toprule
$N$ parity & Critical mode & $\Delta S$ \\
\midrule
Even (self-paired $m = N/2$) & one eigenvalue crosses & $\log 2$ (1 bit) \\
Odd (paired $m, N{-}m$) & two eigenvalues cross & $2\log 2$ (2 bits) \\
\bottomrule
\end{tabular}
\end{center}
\noindent
From $N = 7$ (the flat-plane marginal threshold) through $N = 12$:
exactly $9$ bits (verified numerically).
Through $N = 17$ ($11$ wall-crossings): the discrete part contributes
$17\ln 2 \approx 11.8$ nats ($9.9\%$), the continuous one-loop part
$107.8$ nats ($90.1\%$).
The $L^2$-Morse index $\tau(P_-)$ has denominator~$2$
at each threshold (from the eta invariant shift of $1/2$
due to the zero mode), consistent with the $K$-theoretic framework.

\paragraph{The one-loop correction.}
The one-loop partition function around the polygon saddle is
$Z_{\mathrm{1-loop}} = (\det{}'\Delta_{\mathrm{Lich}})^{-1/2}$,
where $\Delta_{\mathrm{Lich}}$ is the Lichnerowicz operator
on the multi-cone geometry.
By the conformal invariance of the flat-plane
backreaction (Proposition~\ref{II-prop:conformal-invariance}),
the leading correction comes from the spatial topology:
the spectral bound
$|\delta C_1| \le 2.86/(4\pi\lambda_1)$
(where $\lambda_1$ is the spectral gap of the spatial
Laplacian) gives a rigorous bound on the quantum shift
of the stability eigenvalues.
On the Bolza surface ($\lambda_1 \approx 3.84$):
$|\delta C_1| < 0.06$, which does not shift
$N_{\mathrm{crit}}$ away from~$7$ on the flat plane.
On the modular surface ($\lambda_1 \approx 91$):
$|\delta C_1| \lesssim 0.002$, negligible.
The Bolza surface carries the $\{8,3\}$ hyperbolic lattice
on which \citet{Lenggenhager2023} constructed the
hyperbolic Haldane model (Chern insulator with quantised
Bott index $B = \pm 1$ and chiral edge states).
The vortex stability analysis on the same $\{8,3\}$
lattice shares the spectral data ($\lambda_n$ of $\Delta_S$)
with the hyperbolic Haldane model.
The bound $|\delta C_1| < 0.06$ shows that polygon
stability is insensitive to the hyperbolic lattice topology.
\paragraph{Lorentzian continuation and the BTZ horizon.}
The Euclidean analysis above determines the thermal phase structure; the Lorentzian continuation identifies the real-time dynamics.
The Green's function offset $\Phi(\rho) = \log(2\sinh\rho) + b(N)$ admits
an analytic continuation $\rho \to \rho + i\pi/2$ that interchanges
the BTZ exterior and interior:
$\sinh(\rho + i\pi/2) = i\cosh\rho$, so
$\Phi^{\mathrm{int}} = \log(2\cosh\rho) + b(N) + i\pi/2$.
The analytically continued kernel interpolates between curvature
signs:
\begin{center}
\begin{tabular}{lccc}
\toprule
$\mathrm{Im}(\rho)$ & $0$ & $\pi/2$ & $\pi$ \\
\midrule
Curvature sign & $K{=}{-}1$ & $K{=}0$ & $K{=}{+}1$ \\
Continued kernel & $\log(2\sinh\rho)$ & $\log(2\cosh\rho) + i\pi/2$ &
  $\log(2\sinh\rho) + i\pi$ \\
\bottomrule
\end{tabular}
\end{center}
\noindent
(The real parts of the continued kernels reproduce the
pair interactions on spaces of the corresponding curvature
sign; the imaginary parts are phases.)
At each palindromic threshold $\rho^*(N)$: the critical eigenvalue
$\lambda_{m^*}$ crosses zero (\emph{real} side), while the interior
resonance has $\mathrm{Im}(\lambda) = \pi/2$ (universal, independent
of $N$ and~$m$).
The polygon--BTZ transition is a \emph{time--space exchange}:
for $\rho > \rho^*$ the mode oscillates in time
($\omega = \sqrt{\lambda_{m^*}}$); for $\rho < \rho^*$ it grows
spatially ($\kappa = \sqrt{|\lambda_{m^*}|}$);
at $\rho = \rho^*$ the Killing vector degenerates.

Under Lorentzian continuation ($\rho \to \rho + i\varepsilon$),
the partition function acquires a branch cut at the threshold:
$Z_L \sim A(N)\,({\rho - \rho^* + i\varepsilon})^{-1/2}$
with spectral function $\mathrm{Im}[Z_L] \propto |\rho - \rho^*|^{-1/2}$
(a van~Hove singularity).
The exponent $-1/2$ (divergent density at the edge) is the
spectral signature of the integrable Havelock system.
The Dyson circular ensemble (GUE, $\beta = 2$) has the
opposite exponent: $\rho \propto |\lambda - \lambda_{\mathrm{edge}}|^{+1/2}$
(vanishing density, the Wigner semicircle edge).
The sign flip $-1/2 \to +1/2$ is a theorem:
the Havelock gaps satisfy $\lambda_{m^*+\delta} = \delta^2/2$
(even~$N$) or $\delta(\delta{-}1)/2$ (odd~$N$)
(verified numerically for
$N = 7,\ldots,30$; see Code Availability), giving $\rho(\lambda) = (2\lambda)^{-1/2}$
(van~Hove, the integrable exponent \citep{BerryTabor1977}).
The GUE edge exponent $+1/2$ is the Tracy--Widom theorem
\citep{TracyWidom1994}.
The BTZ black hole has GUE level statistics
\citep{AurichSteiner1993}.


%% ====================================================================
\section{The Wheeler--DeWitt equation on the palindromic minisuperspace}
\label{sec:wdw}
The single degree of freedom $\rho$ (geodesic radius of the polygon)
and the potential $V(\rho) = \lambda_{m^*}(\rho)$ (the critical
Havelock eigenvalue) define a minisuperspace Wheeler--DeWitt equation:
\begin{equation}\label{eq:wdw}
  \biggl[-\frac{1}{2c}\,\frac{d^2}{d\rho^2}
  + \Phi(\rho) - \frac{m(N{-}m)}{2}
  \biggr]\Psi = 0,
  \qquad \Phi(\rho) := \log(2\sinh\rho) + b(N),
\end{equation}
with $c = 12\,b(N)$ and $\hbar = 1$
(the semiclassical limit is $c \to \infty$).
On constant-curvature $\mathbf{H}^2$ the Havelock remainder $\delta_m = 0$
identically (\S\ref{sec:three-layer}), so it does not appear.

\emph{Notation.}
We write $\Phi(\rho) = \log(2\sinh\rho) + b(N)$ for the
Green's function offset (the unconstrained radial potential),
reserving $C_1(\xi)$ for the constrained angular stiffness
(eq.~\ref{eq:C1-H2}).
The WDW potential is $V(\rho) = \Phi(\rho) - f(m^*,N)$.

\emph{Logical status of the coefficient $c = 12\,b(N)$.}
The value $b(N) = N(N{+}1)/12 - \ln 2 + (\ln N)/(N{-}1)$ is
computed: it is the unique mode-independent
offset in the Havelock decomposition
$\lambda_m = C_1 - f_m + \delta_m$
(Proposition~\ref{prop:equiv-rr}), proved analytically in
Paper~I via the Gauss product
$\prod_{p=1}^{N-1}\sin(\pi p/N) = N/2^{N-1}$.
Define $c := 12\,b(N)$.
No free parameter is involved.

\emph{Why second-order dynamics from a first-order system.}
The Kirchhoff equations are first-order, but the WDW equation
for the breathing mode~$\rho$ is second-order.
The derivation has two steps:
\begin{enumerate}
\item \emph{Existence (Caldeira--Leggett).}
  The $N{-}2$ angular modes form a bath whose influence functional
  produces a $\dot\rho^2$ term in the Onsager--Machlup path integral
  (Proposition~\ref{prop:wdw-lambda}).
  The Gaussian (CL) approximation gives a coefficient
  $c_{\mathrm{CL}} \sim N^2\rho^{*2}/((N{-}2)\coth^2\rho^*)$
  that is quantitatively wrong (factor~$0.5$--$2$ depending on~$N$),
  because it neglects the higher integrals of the CMS bath.
\item \emph{Exact coefficient (integrability).}
  The CMS angular modes are integrable (Liouville--Arnold):
  the $N{-}2$ angle-action variables decouple, so the one-loop
  partition function $Z_{\mathrm{bath}}(\rho)
  = \prod_{m \ne m^*}|\lambda_m(\rho)|^{-1/2}$ is
  exact at $O(c^0)$: the action-angle factorisation
  gives $Z = (2\pi)^{N-2}\!\int\!\prod dI_m\,e^{-\beta H(I)}$
  (the $\theta$-integrals are exact), and the remaining
  $I$-integrals have anharmonic corrections
  of order~$1/c$ where $c = 12\,b(N) \sim N^2$
  ($1.9\%$ at
  $N = 7$, $0.8\%$ at $N = 11$, decreasing as~$N^{-2}$;
  see Code Availability).
  Integrability prevents mode mixing (the actions~$I_m$
  are independent), so the $1/c$ series converges
  without secular terms.
  The effective action for~$\rho$ is therefore
  $S_{\mathrm{eff}}(\rho)
  = \tfrac{1}{2}\sum_{m \ne m^*}\!\ln|\lambda_m(\rho)|
  + O(1/c)$,
  and the Feynman--Kac theorem maps this to the WDW Hamiltonian
  $H = -1/(2c)\,d^2/d\rho^2 + V(\rho)$ with
  $c = 12\,b(N)$ (the unique mode-independent offset of the
  three-layer decomposition, Proposition~\ref{prop:equiv-rr}).
\end{enumerate}
The CL mechanism proves $c > 0$; the CMS integrability
gives a controlled $1/c$ expansion with computable corrections.

The four-role consistency of $c = 12\,b(N)$ (Section~\ref{sec:intro-gravity}) is preserved.
Each term has a definite origin:
\begin{center}
\small
\begin{tabular}{lll}
\toprule
Term & Expression & Origin \\
\midrule
Quantum kinetic & $-\hbar^2/(2c)\;\Psi''$ &
  Todd offset $c = 12\,b(N)$ (\S\ref{sec:intro-gravity}) \\
Classical potential & $\log(2\sinh\rho)$ &
  $\mathbf{H}^2$ Green's function \\
Vacuum energy & $b(N) \sim c/12$ &
  Orbifold CFT vacuum (Todd class) \\
Matter & $m(N{-}m)/2$ &
  Twist-field dimension (first Chern class) \\
One-loop & $\delta_m$ &
  Havelock remainder (traceless, palindromic) \\
\bottomrule
\end{tabular}
\end{center}

\noindent
\emph{Three regimes.}
\begin{enumerate}
\item \textbf{Classical} ($c \to \infty$):
  the kinetic term vanishes; the constraint reduces to $V(\rho) = 0$,
  i.e., $\lambda_{m^*}(\rho^*) = 0$---the palindromic threshold,
  which is the Einstein equation of the vortex--gravity system.
\item \textbf{Semiclassical} (WKB):
  $\Psi \sim \exp\!\bigl(\pm(i/\hbar)\!\int\!\sqrt{2c\,V}\,d\rho\bigr)$,
  oscillatory in the exterior ($V > 0$, $\sinh$-branch) and evanescent
  in the interior ($V < 0$, $\cosh$-branch).
  The continuation $\rho \to \rho + i\pi/2$ maps between them
  (\S\ref{sec:btz}, Lorentzian paragraph).
\item \textbf{Quantum} (near threshold):
  $\Psi \sim \mathrm{Ai}\bigl((2c\,\coth\rho^*)^{1/3}
  (\rho - \rho^*)\bigr)$,
  smoothing the classical horizon over the width
\begin{equation}\label{eq:airy-width}
  \delta\rho = \frac{1}{(2c\,\coth\rho^*)^{1/3}},
\end{equation}
  where $\coth\rho^* = V'(\rho^*)$ is the surface gravity.
  For large $\rho^*$ ($\coth\rho^* \to 1$):
  $\delta\rho \to (2c)^{-1/3} \approx (2N^2)^{-1/3}$.
  At $c = \infty$: $\delta\rho \to 0$ (sharp classical threshold).
  At finite~$c$: the horizon has quantum width $O(c^{-1/3})$.
\end{enumerate}

\begin{proposition}[Eigenvalue non-crossing at finite $c$]
\label{prop:non-crossing}
For $N \ge 4$ and $c \ge N^2$ (both satisfied for all $N \ge 4$;
the bound requires $\lfloor N/2 \rfloor \ge 2$ for the
maximisation, i.e.\ $N \ge 4$;
$c \ge N^2$ holds since $c = 12\,b(N) > N^2$ for all $N \ge 3$
since $c = 12\,b(N) = N^2 + N - 12\ln 2 + O(\ln N/N) > N^2$
for all $N \ge 3$), the mode-dependent part of the
$O(1/c)$ Virasoro block correction satisfies
\begin{equation}\label{eq:non-crossing}
  \frac{|\delta f_m - \delta f_{m+1}|}{f(m{+}1,N) - f(m,N)}
  = \frac{2m(N{-}m{-}1) + (N{-}1)}{2c}
  \;\le\; \frac{1}{4} - \frac{1}{2N^2}
  \;<\; \frac{1}{4}.
\end{equation}
In particular, no eigenvalue crossing occurs and the critical mode
index $m_{\mathrm{crit}} = \lfloor N/2\rfloor$ is preserved.
\end{proposition}

\begin{proof}
The leading $1/c$ correction to the Casimir is
$\delta f_m = -f_m^2/c = -m^2(N{-}m)^2/(4c)$.
The difference factors as
$\delta f_m - \delta f_{m+1}
= (N{-}2m{-}1)\bigl[2m(N{-}m{-}1)+(N{-}1)\bigr]/(4c)$.
Dividing by the Casimir gap $f(m{+}1)-f(m) = (N{-}2m{-}1)/2$
gives the ratio $R(m) = [2m(N{-}m{-}1)+(N{-}1)]/(2c)$.
The maximum over $1 \le m < N/2$ is at $m = N/2 - 1$, giving
$R_{\max} = (N^2{-}2)/(4N^2) = 1/4 - 1/(2N^2) < 1/4$.
Since $R < 1$ for all adjacent pairs: the residual gap is at least
$3(N{-}2m{-}1)/8 > 0$, and no reordering occurs.
\end{proof}

\begin{remark}[Algebraic field preservation]
\label{rmk:field-preservation}
The Zamolodchikov $c$-recursion gives the correction as a
\emph{rational function} of~$c$, with only
$\mathbb{Z}_N$-invariant channels contributing.
At $c = N^2$: the palindromic coefficients $A(c), B(c)$ remain in
$\mathbb{Q}(\cos(2\pi/N))$---the same number field as at $c = \infty$.
The palindromic form $A\xi^2 + B\xi + A = 0$ is exact at all~$c$
(the conformal inversion $\xi \leftrightarrow 1/\xi$ is a geometric
symmetry of $\mathbf{H}^2$, not a perturbative artefact).
\end{remark}

The palindromic threshold $\xi^*(N)$ is the critical AdS
radius at which the polygon--BTZ transition occurs, and
the algebraic number field $\mathbb{Q}(\sqrt{D})$ of the
threshold is the arithmetic invariant of the phase boundary.

%% ====================================================================

\section{The three-layer decomposition}
\label{sec:three-layer}

\begin{proposition}[Orbifold--Havelock correspondence]
\label{prop:orbifold-havelock}
At leading order in $1/N$, the Havelock eigenvalue on
$\mathbf{H}^2$ matches the $\mathbb{Z}_N$ orbifold CFT
ground-state energy at $c_{\mathrm{orb}} = 12N^2$
(the UV orbifold central charge)
with angular impulse constraint:
\begin{equation}\label{eq:orbifold-havelock}
  \lambda_m = -E_0(m;\,c_{\mathrm{orb}}{=}12N^2) + \mu_L(\xi),
\end{equation}
where $E_0(m) = -N/2 + m(N{-}m)/2$ is the twisted-sector
ground state and $\mu_L(\xi) = C_1(\xi) - N/2$ is the
Lagrange multiplier.
The matching uses $c_{\mathrm{orb}} = 12N^2$
(so that $c/(24N^2) = 1/2$), while the exact Havelock central charge
is $c = 12\,b(N) = N(N{+}1) - 12\ln 2 + \ldots \approx N^2$ at
leading order.
The distinction: $c_{\mathrm{orb}}$ is the UV (free-field) central
charge of the $\mathbb{Z}_N$ orbifold; $c = 12\,b(N)$ is the exact
IR value after regularisation of the vortex self-energy.
The mode-dependent matching (the $m(N{-}m)/2$ term) uses only
$c_{\mathrm{orb}}/(24N^2) = 1/2$; the mode-independent
$O(N)$ corrections (the $\ln 2$ and $(\ln N)/(N{-}1)$ terms in $b(N)$)
enter through~$\mu_L$, not through~$E_0$.
\begin{center}
\small
\begin{tabular}{lll}
\toprule
Symbol & Value & Role \\
\midrule
$c_{\mathrm{orb}} = 12N^2$ & UV orbifold & Mode-dependent matching \\
$c = 12\,b(N) \approx N^2{+}N$ & Exact (IR) & WDW kinetic coefficient \\
$c \approx N^2$ & Leading order & Approximation in estimates \\
\bottomrule
\end{tabular}
\end{center}
\end{proposition}

\begin{proof}
The orbifold ground-state energy
$E_0(m) = -c/(24N) + cm(N{-}m)/(24N^2)$ at $c = c_{\mathrm{orb}} = 12N^2$
(the UV orbifold central charge, which gives $c/(24N^2) = 1/2$ exactly)
yields $E_0 = -N/2 + m(N{-}m)/2$.
Then $-E_0(m) + \mu_L = [N/2 - m(N{-}m)/2] + [C_1(\xi) - N/2]
= C_1(\xi) - m(N{-}m)/2 = \lambda_m$.
\end{proof}

\begin{proposition}[Todd class decomposition of the eigenvalue offset]
\label{prop:equiv-rr}
The $\mathbb{Z}_N$-equivariant eigenvalue decomposition
of the vortex Laplacian on the cusped Poincar\'e disk
(analogous to the equivariant Riemann--Roch theorem) satisfies
\begin{equation}\label{eq:equiv-rr}
  \mathrm{ind}_m = \underbrace{\frac{m(N{-}m)}{2}}_{c_1\text{-term}}
    + \underbrace{b(N)}_{\text{Todd term}}
    + \delta_m,
\end{equation}
where $b(N) = N(N{+}1)/12 - \log 2 + \log N/(N{-}1)$
is independent of both~$m$ and~$\rho$ (\S\ref{sec:intro-gravity},
verified in Paper~I).
The central charge $c = 12\,b(N) = N^2 + O(N)$ is the same
for every mode~$m$, extracted from the Todd correction alone.
\end{proposition}

\begin{proof}
The Havelock eigenvalue satisfies $\lambda_m + f_m = \Phi(\rho) + \delta_m + O(e^{-\rho})$,
where $\Phi(\rho) = \log(2\sinh\rho) + b(N)$ (the radial
dependence is absorbed by the Green's function).
Subtracting $\log(2\sinh\rho)$:
$\mathrm{ind}_m := \lambda_m + f_m - \log(2\sinh\rho) = b(N) + \delta_m$.
Since $\delta_m$ is traceless ($\sum\delta_m = 0$ to $10^{-15}$)
and palindromic ($\delta_m = \delta_{N-m}$), the mode-independent
part $b(N)$ is extracted exactly.
\end{proof}

\noindent
The three-layer decomposition
$\lambda_m = C_1 - f_m + \delta_m$
is the spectral analogue of the equivariant Riemann--Roch theorem:
$f_m$ is the $c_1$-contribution (twist-field dimension),
$b(N)$ is the Todd-class correction (vacuum energy $\approx c/12$),
and $\delta_m$ is the one-loop remainder (the Havelock remainder).
The central charge is $c = 12\,b(N)$,
determined by the index theorem.

\paragraph{Exact Havelock remainder formula.}
\emph{Clarification:} The anomaly $\delta_m$ below is the
\emph{raw} three-layer residual for the \emph{unconstrained} eigenvalue
on the flat plane.
On \emph{constant-curvature} surfaces with the \emph{constrained}
(angular-impulse) eigenvalue, $\delta_m = 0$ identically
(Paper~IV, Theorem~\ref*{IV-thm:zero-anomaly}).
On constant-curvature surfaces, $\delta_m = 0$ is a consistency
check (the decomposition is exact where it should be).
The non-trivial content is on compact quotients
$\Gamma\backslash\mathbf{H}^2$: the Selberg trace formula gives
$\delta_m = \delta C_1$ (mode-independent) with
$|\delta C_1| \le 0.020$ on the Bolza surface
(\S\ref{sec:btz}).
The three-layer decomposition is therefore exact up to a
uniform shift even on non-constant-curvature surfaces.
The nonzero formula below applies to the flat-plane or
non-constant-curvature case.

\noindent
The anomaly $\delta_m$ decomposes exactly into character sums
on~$\mathbb{Z}/N\mathbb{Z}$:
\begin{equation}\label{eq:weyl-exact}
  \delta_m = L(m,N) - \tfrac{1}{4}\,S_2(m,N) + c(N),
\end{equation}
where $L(m,N) = \sum_{p=1}^{N-1}\cot(\pi p/N)\cos(2\pi pm/N)$
is the cotangent sum,
$S_2(m,N) = \sum_{p=1}^{N-1}\csc^2(\pi p/N)\cos(2\pi pm/N)$
is the cosecant-squared sum, and $c(N)$ is the unique constant
enforcing $\sum_{m=1}^{N-1}\delta_m = 0$ (tracelessness).
The coefficients $a = 1$, $b = -1/4$ are exact and universal
($R^2 = 1.000000$ for all tested $N = 4,\ldots,24$).
This identifies the anomaly as the lattice correction
(orbifold fixed-point contribution in the KK picture)
to the continuum Casimir.

The value $c = 12\,b(N) \approx N^2$ makes five independently
testable predictions beyond the Casimir:
(i)~the exact offset $b(N) = N(N{+}1)/12 - \log 2 + \log N/(N{-}1)$
(Paper~I);
(ii)~the mean aliasing $\langle D\rangle = N(N{+}1)/12$ (exact);
(iii)~the growth law $\rho^*/f_{\mathrm{crit}} \to 1/3$ (Paper~I);
(iv)~the quartic coupling $Q/f_{\mathrm{crit}}^2 = N/48$ (exact for even~$N$);
(v)~the mode-independent Laplacian $\Delta_{\mathbf{H}^2} h = -1/2$.
All five follow from $c = 12\,b(N) \approx N^2$ through the $B_2$-tower:
the Casimir $m(N{-}m)/2$ is $O(c)$; the Havelock remainder $\delta_m$
is $O(c^0)$ (topological); the Virasoro block correction is
$O(1/c) \approx 10^{-5}$ (\S\ref{sec:holographic}).

The tower extends to a \emph{Magri bihamiltonian hierarchy}.
As noted in \S\ref{sec:graviton}, the Thomson vortex system has the same linearised
spectrum as the CMS model at coupling $g = 1$.
The CMS model is bihamiltonian
\citep{OlshanetskyPerelomov1976}:
the Kirchhoff bracket $P_0$ (mode-independent) and the
Casimir bracket $P_1$ (eigenvalue $f_m$ on mode~$m$) are
compatible, with recursion operator $R = P_1 P_0^{-1}$
acting as multiplication by $f_m = m(N{-}m)/2$.
By the Magri theorem, the $k$-th conserved quantity $I_k$
has mode eigenvalue $f_m^k$,
with physical derivatives
$d_{2k} = \alpha_k N f_m^k$
($\alpha_1 = 1$, $\alpha_2 = 4$, $\alpha_3 = 64$,
verified numerically for $N \le 16$).
The quartic coupling $Q = (N/48)f^2$ lies inside the hierarchy
($H_4^{\mathrm{diag}} \propto I_2$), guaranteeing
$\{I_j, H_4\} = 0$ by integrability.
The three-layer decomposition $\lambda_m = C_1 - f_m + \delta_m$
is the truncation of this hierarchy at level~$1$.

On arithmetic surfaces (Bolza, modular, etc.), the eigenvalue
acquires a topological correction that is invisible to the
confining potential $C_1(\xi)$:
\begin{equation}\label{eq:three-layer}
  \lambda_m = \underbrace{C_1(\xi)}_{\text{Ricci scalar}}
    - \underbrace{\frac{m(N{-}m)}{2}}_{\text{Casimir}}
    + \underbrace{\delta_m(N)}_{\text{Havelock remainder}}.
\end{equation}
Each layer has a distinct character:
\begin{enumerate}
\item $C_1(\xi)$: continuous, surface-dependent, mode-independent.
  Scales approximately as $C_1 \propto \rho$ over the range
  $\rho \in [0.01, 10]$ (with $R^2 = 0.98$,
  where $\rho$ is the geodesic radius
  and $\xi = \tanh^2(\rho/2) \in (0,1)$);
  the exact formula is eq.~\eqref{eq:C1-H2}.
\item $f(m,N) = m(N{-}m)/2$: discrete, universal (surface-independent).
  The $\mathbb{Z}_N$ Casimir from the Havelock identity.
\item $\delta_m(N)$: topological, UV-finite, $\xi$-independent
  ($\text{variation} < 10^{-13}$), and exactly traceless
  ($\sum_{m=1}^{N-1} \delta_m = 0$ to $10^{-15}$).
  Carries all number-theoretic content: Hecke eigenvalues,
  palindromic structure, Langlands data.
\end{enumerate}
By structural analogy with the decomposition of the Riemann
curvature tensor: $C_1$ plays the role of the Ricci scalar,
$f(m,N)$ the traceful part, and $\delta_m$ the traceless Weyl-like
anomaly.
In AdS$_3$/CFT$_2$, the Brown--Henneaux scaling
$C_1 \propto \log(\xi)$ says that
the Weyl-like anomaly $\delta_m$ is the topological correction
that carries the arithmetic of the underlying surface, contributing
corrections of order $\|\delta\|/C_1 \sim 2\%$ at typical parameter values.

The Weyl norm scales as
$\|\delta\|^2 \sim N^{4.2}(\log N)^{4.1}$ for $N \ge 40$
(not a pure power law).
The triangle ($N = 3$) is \emph{anomaly-free}:
$\delta_m = 0$ exactly for all Green's functions and curvatures,
because $\cos(2\pi/3) = \cos(4\pi/3)$ forces $G_1 = G_2$.
Under Casimir-conserving mergers $N_1 + N_2 \to N_3$
($N_3^3 {-} N_3 = N_1^3 {-} N_1 + N_2^3 {-} N_2$):
the unique anomaly-\emph{destructive} reaction is
$4{+}4 \to 5$ ($88\%$ Weyl loss); all other mergers create
Weyl content (gains $46\%$ to $849\%$).
The Weyl entropy per degree of freedom $s(N) = \log\|\delta\|^2/\sum f$
peaks at $N \approx 9$--$10$, coinciding with the observed
planetary polygon range ($N = 5, 6, 8$).

The anomaly $\delta_m$ has two independent origins, which the
parameter space $(\Delta, \xi)$ separates cleanly:
\begin{enumerate}
\item \emph{Interaction anomaly} ($\Delta > 0$ on $\mathbb{R}^2$):
  deviation of $V = r^{-2\Delta}$ from $-\log r$.
  Large ($\|\delta\|^2 = O(1$--$100)$) but arithmetically trivial.
\item \emph{Geometric anomaly} ($\Delta = 0$ on $\mathbf{H}^2$):
  deviation of the hyperbolic Green's function from the flat one.
  Small ($\|\delta\|^2 = O(0.01$--$100)$) but carries all
  number theory: palindromic polynomials, Hecke eigenvalues,
  Langlands data.
\end{enumerate}
The RG flow ($\Delta \to 0$) strips away the large, trivial
interaction anomaly, leaving behind the small geometric anomaly
as the residue that survives at the conformal fixed point.
The number theory is not emergent from the flow; it is a property
of the conformal fixed point, activated by curvature.
The three-layer structure has a symplectic origin: the Kirillov
form $\omega_{\mathrm{KR}}$ (mode-independent) and the Weil--Petersson
form $\omega_{\mathrm{WP}}$ (mode-dependent) are \emph{not} proportional;
their discrepancy $\omega_{\mathrm{WP}} - \alpha\,\omega_{\mathrm{KR}}
\propto f(m,N)$ is the Casimir (Layer~2), and the residual beyond
this linear relation is the Havelock remainder (Layer~3).

The $B_2$ organising principle --- connecting the stability polynomial
$6B_2(x) = P_2(2x{-}1)$ to Gauss--Legendre quadrature, the
Euler--Maclaurin duality, and the $\csc^2$ kernel universality
(Calogero--Moser--Sutherland, Dyson, Haldane--Shastry, Laughlin,
Selberg, $W_N$, AGT/Nekrasov) ---
is developed in Paper~I, \S\ref*{I-sec:b2-tower}.

%% ====================================================================

\section{Higher-dimensional extension: the CFT$_3$/AdS$_4$ Casimir}
\label{sec:cft3}

The two-dimensional results extend to conformal field theory in
any dimension~$d$ via the power-law interaction
$V(x,y) = |x - y|^{-2\Delta}$, where $\Delta$ is the conformal
dimension. For $\Delta \to 0$: $V \to -\log|x-y|$ (the 2D case).
For $\Delta = 1$: the CFT$_3$ two-point function on $S^2$
(the boundary of AdS$_4$).

\begin{proposition}[Generalised Havelock Casimir]
\label{prop:generalised-casimir}
For $N$ equally-spaced points on a ring on $S^2$,
interacting via $V = |x - y|^{-2\Delta}$,
the angular Hessian eigenvalues factorise as
$\mu_m = g(\theta_0, \Delta) \cdot h(m, N, \Delta)$,
where the universal Casimir is
\begin{equation}\label{eq:gen-casimir}
  h(m, N, \Delta) = \sum_{p=1}^{N-1}
    \frac{1 + 2\Delta\cos^2(\pi p/N)}
         {\sin^{2\Delta+2}(\pi p/N)}\,
    \bigl(1 - \cos\tfrac{2\pi mp}{N}\bigr).
\end{equation}
For $\Delta = 0$: $h(m, N, 0) = 2m(N{-}m)$.
The ratio $h(m)/h(1)$ is independent of $\theta_0$, confirming
that the Havelock decomposition extends to all~$\Delta$.
\end{proposition}

The constrained Hessian on the flat plane with $V = |x - y|^{-2\Delta}$
gives an $N$-dependent stability boundary $N_{\mathrm{crit}}(\Delta)$:

\begin{center}
\begin{tabular}{lcccccc}
\toprule
$\Delta$ & $0$ & $0.5$ & $1$ & $2$ & $5$ & $10$ \\
\midrule
$N_{\mathrm{crit}}$ & $7$ & $5$ & $5$ & $5$ & $4$ & $4$ \\
\bottomrule
\end{tabular}
\end{center}

The classical $N_{\mathrm{crit}} = 7$ is an isolated phenomenon
of the logarithmic interaction. For the CFT$_3$ value $\Delta = 1$:
$N_{\mathrm{crit}} = 5$.

Under the Green's function identification $\Delta = (d{-}2)/2$,
the table acquires a dimensional interpretation:
\begin{center}
\begin{tabular}{lccccc}
\toprule
Dimension $d$ & $2$ & $3$ & $4$ & $5$ & $\ge 6$ \\
$\Delta = (d{-}2)/2$ & $0$ & $1/2$ & $1$ & $3/2$ & $\ge 2$ \\
\midrule
$N_{\mathrm{crit}}$ & $7$ & $5$ & $5$ & $5$ & $4$--$5$ \\
\bottomrule
\end{tabular}
\end{center}
\noindent
The hexagon ($N = 6$, Saturn) is a \emph{uniquely two-dimensional}
phenomenon: it is marginally stable at $d = 2$ ($\lambda_{\min} = 1/2$)
but unstable for $d \ge 3$ ($\lambda_{\min} = -0.02$ at $d = 3$;
$N_{\mathrm{crit}}(\Delta{=}0.5) = 5 < 6$; see Code Availability).
The pentagon ($N = 5$) survives in all tested dimensions $d \le 10$.

\begin{proposition}[Soft Havelock decomposition in dimension $d$]
\label{prop:soft-havelock}
For the interaction $V = r^{-(d-2)}$ ($d \ge 3$) on the $N$-gon ring,
the angular Casimir decomposes as
\begin{equation}\label{eq:soft-havelock}
  \sigma_m^{(d)} = a(N,d)\,\frac{m(N{-}m)}{2} + R_m(N,d),
\end{equation}
where $a(N,d)$ is the leading Casimir coefficient (determined by
least-squares fit) and $R_m$ is the residual.
For all tested $N = 5,\ldots,20$ and $d = 3,4,5$: the
residual satisfies $|R_m - R_{m'}| < a\,|f_m - f_{m'}|$ for every
adjacent pair, so the \emph{mode ordering is preserved}.
The critical mode is always $m = \lfloor N/2 \rfloor$ (same as $d=2$),
and the stability threshold has the same palindromic form as the
two-dimensional case, with coefficients that depend on~$d$.
\end{proposition}

\noindent
For $d = 2$ ($s = 0$): $a = 1$ exactly and $R_m = 0$ (the Havelock
identity). This is the \emph{unique} dimension where the Casimir is
an exact polynomial---because $\csc^2(\pi p/N)$ has a closed discrete
Fourier transform ($\sum\csc^2\cdot(1{-}\cos) = 2m(N{-}m)$), while
$\csc^s$ for $s \ne 2$ does not.
The exact algebraic structure of Paper~I (palindromic polynomials,
algebraic number fields, modular forms) is specific to $d = 2$,
a consequence of the uniqueness of the $\csc^2$ weight function
among power-law potentials.
However, the \emph{structural skeleton} survives in all dimensions:

\begin{enumerate}
\item \textbf{Mode ordering.}
  The critical mode is $m = \lfloor N/2\rfloor$ in every $d$
  (Proposition~\ref{prop:soft-havelock}).
  The residual $R_m$ is too small to reorder eigenvalues.

\item \textbf{Palindromic symmetry.}
  The symmetry $\sigma_m = \sigma_{N-m}$ holds exactly for all~$d$
  (from the discrete cosine: $\cos(2\pi pm/N) = \cos(2\pi p(N{-}m)/N)$).
  On curved spaces, conformal inversion $\xi \leftrightarrow 1/\xi$ is
  geometric (a property of the Poincar\'e ball, not of the interaction),
  so the threshold equation retains palindromic form.

\item \textbf{Finite-$c$ algebraic structure.}
  The $O(1/c)$ Virasoro block correction
  $\delta h_m = -f_m^2/c$ is a \emph{polynomial in the Casimir}
  (equation~\ref{eq:block-1-over-c}),
  preserving the algebraic structure at every order.

\item \textbf{Havelock remainder shape.}
  The normalised anomaly profile $\delta_m/\|\delta\|$
  has inner products $> 0.98$ between $\Delta = 0$ and
  $\Delta = 2$ for $N \le 6$: the number-theoretic fingerprint
  transfers to $d = 3$ with $< 2\%$ distortion.
\end{enumerate}
The $d = 2$ case is where the algebra is exact;
the $d \ge 3$ case is where it is \emph{approximate but structurally
preserved}.  The mechanism by which the logarithmic structure
survives in $3{+}1$D is the soft Havelock decomposition:
the Casimir $m(N{-}m)/2$ is the dominant term at every~$d$, with
the interaction-specific residual $R_m$ as a controlled perturbation.

\begin{remark}[Transcendence of the transition dimensions]
\label{rmk:transcendence}
The critical value $\Delta^*(N)$ where $N_{\mathrm{crit}}$ drops
below~$N$ is expected to be \emph{transcendental} for $N \ge 5$.
The threshold equation involves $\sin^{-2\Delta}(\pi p/N)$ for
different~$p$ with $\mathbb{Q}$-linearly independent logarithms.
For $N = 5$: the threshold equation involves
$\sin^{-2\Delta}(\pi/5)$, an algebraic base raised to the
power~$\Delta$.
If $\Delta^*$ is algebraic and irrational, the
Gelfond--Schneider theorem makes the individual terms
transcendental, and Baker's theorem on linear forms in
logarithms would be needed to show no cancellation occurs.
A rigorous proof requires Baker-type estimates.
This contrasts with the palindromic thresholds $\xi^*(N)$
on $\mathbb{H}^2$, which lie in explicit algebraic number fields.
\end{remark}

For $N = 5$, the pentagon identity
$\sin(2\pi/5)/\sin(\pi/5) = \varphi$ (the golden ratio)
gives the eigenvalue as
$\lambda_2(\Delta) = P_0(\Delta) + P_1(\Delta)\,\varphi^{-(2\Delta+2)}$,
where $P_0, P_1$ are algebraic functions of $\Delta$.
The transition occurs at $\varphi^{2\Delta^*+2} = P_1/({-P_0})
\approx 25.8$, giving $\Delta^* \approx 2.378$.

\begin{theorem}[RG monotonicity of the Casimir ratio]
\label{thm:rg-monotonicity}
For any $N \ge 4$ and $m_{\mathrm{crit}} = \lfloor N/2 \rfloor$,
the generalised Casimir ratio
$R(\Delta) = h(m_{\mathrm{crit}}, N, \Delta) / h(1, N, \Delta)$
is strictly increasing on $[0, \infty)$, with exact limits
\begin{equation}\label{eq:rg-limits}
  R(0) = \frac{m_{\mathrm{crit}}(N{-}m_{\mathrm{crit}})}{N{-}1},
  \qquad
  R(\infty) = \biggl[\frac{\sin(\pi m_{\mathrm{crit}}/N)}
                          {\sin(\pi/N)}\biggr]^2.
\end{equation}
\end{theorem}

\begin{proof}
By the Chebyshev identity, the mode amplification factor
$f(p) = (1{-}\cos(2\pi m_{\mathrm{crit}} p/N))/(1{-}\cos(2\pi p/N))
= [\sin(\pi m_{\mathrm{crit}} p/N)/\sin(\pi p/N)]^2$
achieves its global maximum at $p = 1$ (and its palindromic twin
$p = N{-}1$):

\emph{Even $N = 2k$:}
$\sin(\pi kp/N) = \sin(\pi p/2) \in \{0, \pm 1\}$, so
$f(p) = 1/\sin^2(\pi p/N)$ for odd~$p$ and $f(p) = 0$ for even~$p$.
Since $\sin(\pi p/N) \ge \sin(\pi/N)$ for $p \ge 2$: $f(p) \le f(1)$.

\emph{Odd $N = 2k{+}1$:}
The identity $\sin(k\pi p/N) = \sin(\pi p/2 - \pi p/(2N))$
(using $k/N = 1/2 - 1/(2N)$) gives explicit formulas:
$f(p) = 1/(4\sin^2(\pi p/(2N)))$ for odd~$p$, and
$f(p) = 1/(4\cos^2(\pi p/(2N)))$ for even~$p$.
Since $f(1) = 1/(4\sin^2(\pi/(2N)))$:
for odd $p \ge 3$, $\sin(\pi p/(2N)) \ge \sin(\pi/(2N))$
(monotonicity of $\sin$ on $(0,\pi/2)$), so $f(p) \le f(1)$;
for even $p \le N{-}1$, $\cos(\pi p/(2N)) \ge \sin(\pi/(2N))$
(since $\pi p/(2N) + \pi/(2N) \le \pi/2$ iff $p \le N{-}1$),
so $f(p) \le f(1)$.

The normalised weight $\mu_p(\Delta) \propto \tilde{w}(p,\Delta)\,c_1(p)$
concentrates on $\{1, N{-}1\}$ as $\Delta \to \infty$, because
$\tilde{w}(p)/\tilde{w}(1) = O([\sin(\pi/N)/\sin(\pi p/N)]^{2\Delta})
\to 0$ for $p \notin \{1, N{-}1\}$.
Since $R(\Delta) = \mathbb{E}_\mu[f]$ and $\mu$ concentrates at the
global maximum of~$f$:
$dR/d\Delta = \mathrm{Cov}_\mu(d\log\tilde{w}/d\Delta,\, f) > 0$
(both $f$ and $g = d\log\tilde{w}/d\Delta = -2\log\sin(\pi p/N)$
are maximised at $p = 1$ and monotonically ordered on
$\{1,\ldots,\lfloor N/2\rfloor\}$; the Chebyshev sum inequality
gives $\mathrm{Cov}_\mu(g, f) > 0$).
\end{proof}

\begin{corollary}[Stability monotonicity]
\label{cor:c-theorem}
$N_{\mathrm{crit}}(\Delta)$ is non-increasing in~$\Delta$: once a
polygon type loses stability, it never regains it.
The quantity $\log N_{\mathrm{crit}}(\Delta)$ decreases from
$\log 7$ ($\Delta = 0$, the Havelock regime) to
$\log 3$ ($\Delta \to \infty$, only triangles survive).
This is reminiscent of the Zamolodchikov c-theorem,
though $N_{\mathrm{crit}}$ is a step function (not smooth)
and $\Delta$ is not an RG parameter.
\end{corollary}

%% ====================================================================
\section{Triangle universality and dynamical triangulation}
\label{sec:triangle}

The equilateral triangle ($N = 3$) is the \emph{unique} regular
polygon that is dynamically stable for all interaction ranges
and quantum deformation levels.

\begin{theorem}[Triangle universal stability]
\label{thm:triangle-universality}
The equilateral triangle ($N = 3$) is the unique regular polygon
that is stable on the flat plane for all $\Delta \ge 0$,
all TQFT levels $k \ge 3$, and whose tessellation
$\{3,6\}$ is the unique stable regular lattice for all $\Delta > 0$.
\end{theorem}

\begin{proof}[Proof sketch]
Three independent arguments:
\begin{enumerate}
\item \emph{Equal chords.} The triangle is the unique regular polygon
  where all chord distances are equal
  ($\sin(\pi/3) = \sin(2\pi/3) = \sqrt{3}/2$).
  No range-dependent competition between nearest and next-nearest
  neighbors $\implies$ stability is $\Delta$-independent.

\item \emph{Single mode.} The triangle has one nontrivial $\mathbb{Z}_3$
  mode ($m = 1$, with $m = 2$ its palindromic conjugate).
  The Havelock eigenvalue $\lambda_1 = (N{-}1) - 1 = 1 > 0$.
  Mode competition requires at least two independent modes ($N \ge 4$).

\item \emph{Zero Havelock remainder.}
  The anomaly $\delta_m$ is traceless: $\sum \delta_m = 0$.
  With one mode: $\delta_1 = 0$ identically.
  Verified: anomaly norm $= 2.2 \times 10^{-16}$ (machine zero) on
  the square torus.
\end{enumerate}
For the tessellation claim: the dynamical matrix (phonon spectrum)
of the triangular lattice under $V = r^{-2\Delta}$ has all eigenvalues
positive for \emph{every} $\Delta > 0$ (verified by lattice sums with
convergence study and Richardson extrapolation).
The minimum phonon frequency is at the $K$-point of the Brillouin zone.
The square lattice (coordination~$4$) has a permanent $X$-point
instability (stripe mode: $\omega^2(X) < 0$ for all~$\Delta$);
the hexagonal lattice (coordination~$3$) has a permanent $K$-point
instability. The triangular lattice is the \emph{only} regular
lattice that is ever stable, at any interaction range.
\end{proof}

\begin{remark}[Dynamical triangulation]
\label{rmk:cdt}
In causal dynamical triangulations (CDT) and Regge calculus,
spacetime is discretised into simplices (triangles in~2D).
This triangulation is imposed by hand as a regularisation.
Theorem~\ref{thm:triangle-universality} suggests it may be imposed
by physics: among regular lattices, the triangular lattice is the
unique stable configuration, so the gravitational path integral
$Z = \sum_{\text{tilings}} e^{-S}$ may be dominated by
triangulated configurations---the square and hexagonal alternatives
are saddle points with negative phonon modes at all~$\Delta$.
The extension to irregular triangulations (as used in CDT)
remains open.
\end{remark}

%% ====================================================================

\section{The polygon partition function}
\label{sec:gravity-partition}

The one-loop partition function of the vortex system organises
naturally by algebraic number field.

\begin{proposition}[Polygon partition function]
\label{prop:gravity-partition}
The saddle-point partition function for the vortex-polygon system
(a discrete sum over polygon number~$N$, each contributing
a single conical geometry---not a path integral over metrics)
on $\mathbf{H}^2$ is
\begin{equation}\label{eq:Z-gravity}
  Z_{\mathrm{grav}} = \sum_{N=3}^{\infty}
  e^{-b(N)}\,Z_{\mathrm{frozen}}(N),
\end{equation}
where $b(N) = N(N{+}1)/12 - \log 2 + \log N/(N{-}1)$ is the
gravitational action and $Z_{\mathrm{frozen}} = 2^{3(N-2)/4}/(N{-}3)!!$
is the one-loop determinant (even~$N$; the odd-$N$ case is computed
from Casimir gaps).
The series converges super-exponentially: $e^{-b(N)} \sim e^{-N^2/12}$.
Numerically, $Z_{\mathrm{grav}} = 1.3304$ and
$N = 3,\ldots,8$ capture $99.99\%$ of the total.
\end{proposition}

\begin{proof}
Convergence: $e^{-b(N)} = O(e^{-N^2/12})$ decays faster than any
polynomial, while $Z_{\mathrm{frozen}} = O(N^{O(N)})$ grows
sub-exponentially; the product vanishes.
The numerical value $Z = 1.3304$ is computed by direct summation
to $N = 20$ (the tail $N > 20$ contributes $< 10^{-20}$).
\end{proof}

The partition function decomposes by \emph{trace field}---the
algebraic number field generated by $\cos(2\pi/N)$:
\begin{center}
\small
\begin{tabular}{lcr}
\toprule
Trace field & $N$ values & Fraction of $Z$ \\
\midrule
$\mathbb{Q}$ (rational) & $3, 4, 6$ & $90.98\%$ \\
$\mathbb{Q}(\sqrt{5})$ (golden) & $5, 10$ & $8.26\%$ \\
$\mathbb{Q}(\sqrt{2})$ (silver) & $8$ & $0.42\%$ \\
$\mathbb{Q}(\cos 2\pi/7)$ (cubic) & $7$ & $0.34\%$ \\
All higher fields & $\ge 9$ & $< 0.01\%$ \\
\bottomrule
\end{tabular}
\end{center}
\noindent
The rational field dominates; the golden ratio field gives the
leading correction; the entire algebraic tower is exponentially
suppressed by $e^{-N^2/12}$.

The Baum-Connes assembly map $\mu\colon K_0^{\mathbb{Z}_N}(\mathbf{H}^2)
\to K_0(C^*_r(\mathbb{Z}_N))$ sends the equivariant index
$\mathrm{ind}_m = f(m,N) + b(N)$ to the K-theory class with trace
$c = 12\,b(N)$.
(The three Kasparov module conditions for the bounded triple
$\alpha_N = (E, \pi, F_\xi)$ --- self-adjointness, $F^2{-}1$ compact,
and $[\pi(a), F]$ compact --- are verified explicitly in
Paper~I, \S\ref*{I-sec:kk-programme}.)  This is mode-independent (the same~$c$ for
all~$m$) and approximately integer: $c = N^2 + O(N)$ for all tested
$N = 3,\ldots,24$.  The discrepancy $c - N^2 \approx N + O(1)$
is the boundary self-energy ($N/12$) plus regularisation terms.

%% ====================================================================
\section{From polygon stability to the Einstein equations}
\label{sec:lichnerowicz}

The Havelock eigenvalue structure characterises the vacuum
Einstein equations.
In $d = 2$ ($2{+}1$D), the three conditions below follow
from the vortex dynamics (two by theorem, one by the Onsager
principle via the vortex-gravity isomorphism):
the Einstein equations emerge as the entropy-maximising
geometry, not as an additional input.
In $d \ge 3$ ($3{+}1$D and higher), condition~(H1)
becomes \emph{equivalent} to the vacuum Einstein equations
(see Remark~\ref{rmk:higher-d}); the theorem then provides a new
characterisation of Einstein manifolds via polygon stability, rather
than a one-directional derivation.
The full $4$D Einstein equations are derived independently in
Paper~IV through the boundary Virasoro mechanism (the emergent
spin-$2$ graviton of Theorem~\ref*{IV-thm:graviton-4d}).
We state the theorem first, then collect the proof.

\begin{remark}[Logical status]
\label{rmk:logical-status}
(H1) and~(H3) are mathematical theorems
($\csc^2$ uniqueness and Hadamard parametrix).
(H2a) is a theorem (strict concavity $+$ Jensen).
(H2b) is the Onsager variational principle
\citep{Onsager1949,CagliotiFLMP1992,Kiessling1993}
applied via the vortex-gravity isomorphism.
The saddle-point approximation at finite~$N$ is
controlled by $c = 12\,b(N)$ (Laplace method,
$O(1/c) \approx 2\%$ at $N = 7$;
\citealt{Erdelyi1956}, \S2.4).
\end{remark}

\begin{theorem}[Polygon entropy maximisation implies constant curvature]
\label{thm:polygon-einstein}
Let $(S, g)$ be a two-dimensional Riemannian manifold.
Place $N$ equal point vortices in a regular polygon ring
of geodesic radius~$\varepsilon$ centred at~$x \in S$,
interacting via the scalar Green's function of~$(S, g)$.
If:
\begin{enumerate}
\item[\textup{(H1)}]
  \emph{Havelock structure} :
  the angular Hessian eigenvalues
  have the form $\lambda_m = C_1(x,\varepsilon) - m(N{-}m)/2$
  with $C_1$ independent of~$m$.
  Proved by the $\csc^2$ uniqueness
  theorem (Paper~I): the Laplacian Green's function on any
  surface has the $\csc^2$ angular kernel, whose $\mathbb{Z}_N$
  Fourier transform is $m(N{-}m)/2$.
\item[\textup{(H2)}]
  \emph{Entropy maximisation}
  \textbf{[characterisation theorem $+$ Onsager selection]}:
  the polygon entropy functional
  $S_{\mathrm{poly}}[g] = \int S(x)\,dA$, where
  $S(x) = \tfrac{1}{2}\sum_{m \ne m^*}\log\lambda_m(x)$,
  is \emph{strictly concave} in~$R(x)$
  ($d^2S/dC_1^2 < 0$; \S\ref{sec:einstein-derivation}).
  Subject to fixed total energy
  $\int V(\rho(x))\,dA = E$,
  the unique maximum is $R(x) = \mathrm{const}$
  (Jensen's inequality).
  The Onsager selection principle
  \citep{Onsager1949,CagliotiFLMP1992}---that the
  microcanonical ensemble at negative temperature is
  dominated by the entropy-maximising macrostate---selects
  this unique maximum.
\item[\textup{(H3)}]
  \emph{Hadamard universality} :
  condition~\textup{(H1)} holds
  at every point~$x \in S$.
  The Hadamard parametrix
  $G(x,y) = -(1/2\pi)\log d(x,y) + \gamma_S(x) + O(d^2)$
  is universal on any Riemannian surface.
\end{enumerate}
Then the scalar curvature of~$S$ is constant:
$R(x) = R_0$ for all~$x$.
In $2{+}1$D, this gives the vacuum Einstein equations
$R_{\mu\nu} = \Lambda\,g_{\mu\nu}$ (Schur's lemma;
the Weyl tensor vanishes identically in three dimensions).
\end{theorem}

\noindent
The cosmological constant $\Lambda$ is determined by a
separate dynamical argument (the WDW equation,
Proposition~\ref{prop:wdw-lambda} below), not by
the theorem above.
The theorem establishes the \emph{form} of the Einstein
equations; the WDW equation determines their
\emph{content} ($\Lambda$).

\begin{corollary}[Onsager contraction --- (H2b) as theorem]
\label{cor:onsager-contraction}
Define the Onsager iteration $T\colon g \mapsto g'$ on
compact metrics by:
(i)~compute the Onsager mean-field \emph{density}
$\rho(x)$ of $N$ equal-circulation vortices on~$(S,g)$
at negative temperature
(exists by the Onsager--Regge Lemma);
(ii)~map $\rho(x)$ to a curvature field $R'(x) = \delta \times \rho(x)$
via the vortex-gravity isomorphism (equal circulations
$\Rightarrow$ equal deficit angles $\Rightarrow$
$R \propto \rho$);
(iii)~define~$g' = T(g)$ as the metric with curvature~$R'$.

\emph{The iteration operates on the scalar field $R(x)$,
not on polygon geometry.}
The $N$ vortices are individual point masses whose positions
are drawn from the mean-field density~$\rho(x)$;
they do \emph{not} form a regular $N$-gon
(except on a constant-curvature surface, where $\rho = \mathrm{const}$
and the entropy-maximising configuration is the polygon).
The energy landscape $V(x)$ that determines~$\rho$
depends on the local curvature~$R(x)$ through the
Hadamard parametrix $C_1(x) = (N{-}1)[1 - R(x)\varepsilon^2/4 + \cdots]$
---a property of the Green's function of~$(S,g)$ at the point~$x$,
not of any global vortex configuration.

Then the curvature variance contracts:
\begin{equation}\label{eq:onsager-contraction}
  \mathrm{Var}(R(T(g)))
  \le |\beta_{\mathrm{eff}}|^2 R_0^2\,\mathrm{Var}(R(g)),
\end{equation}
where $\beta_{\mathrm{eff}} = 1/f(m^*,N) = 1/C_1$ is the inverse
Havelock stiffness (the effective temperature of the angular bath
is set by the stiffest mode, not by the total polygon energy),
$R_0 = 2\pi\chi/\mathrm{Vol}$ is the mean curvature
($|R_0| = 1$ for the Bolza surface), and
$V_1, V_2$ are defined in the proof below.
At~$N = 7$: $|\beta_{\mathrm{eff}}| R_0 = 1/f(3,7) = 1/6 \approx 0.17$,
giving $97\%$ variance reduction per step.
The contraction improves with~$N$: $|\beta_{\mathrm{eff}}| R_0
= 1/8$ at $N = 8$, $1/15$ at $N = 11$.
\end{corollary}

\begin{proof}
The Onsager equilibrium distributes vortices with mean-field
density $\rho(x) \propto \exp(\beta_{\mathrm{eff}}\,V(x))$,
where $V(x) = V(C_1(R(x)))$ is the BO potential and
$\beta_{\mathrm{eff}} = 1/C_1$ is the inverse stiffness
(the angular bath's response to curvature perturbations is
controlled by $C_1 = f(m^*,N)$, not by the total energy).
For a perturbation~$\delta R$ around mean~$R_0$:
$\rho(x) = \rho_0\exp\!\bigl(\beta_{\mathrm{eff}}
(V_1\,\delta R + \tfrac{1}{2}V_2\,\delta R^2)\bigr)$
where $V_1 = \coth(\rho_*)$ and
$V_2 = -\operatorname{csch}^2(\rho_*)$.
The $N$ vortices have \emph{equal} circulations~$\kappa$,
so each deficit angle $\delta_k = 8\pi G\kappa$ is identical.
The new curvature field is therefore
$R'(x) = \delta \times \rho(x)$
(curvature $=$ deficit angle $\times$ density;
no triangulation is needed).
The contraction operates on the density field:
\[
  \mathrm{Var}(R') \;\le\;
  |\beta_{\mathrm{eff}}|^2 R_0^2\,\mathrm{Var}(R)
  \times \bigl(1 + 2|V_2/V_1|\sqrt{\mathrm{Var}(R)}\bigr)^2.
\]
At~$N = 7$: $|\beta_{\mathrm{eff}}| R_0 = 1/6$,
$|V_2/V_1| = |\operatorname{csch}^2(\rho^*)/\coth(\rho^*)| = 0.12$.
The nonlinear factor
$(1 + 0.24\sqrt{\mathrm{Var}})^2 \le 1.75$
for $\mathrm{Var}(R) \le 2$, giving total contraction
$(1/6)^2 \times 1.75 = 0.049 < 1$.
The condition $|\beta_{\mathrm{eff}}| R_0 < 1$ is
$f(m^*,N) > |R_0|$, which holds for all $N \ge 4$
(since $f(m^*,N) \ge 2$ and $|R_0| = 1$ on the Bolza surface).
The contraction holds rigorously for all
initial $\mathrm{Var}(R) \le 2$.
The iteration $T^n(g_0) \to g^*$ converges geometrically;
by~(H2a), $g^*$ has $R = \mathrm{const}$.
\end{proof}

\begin{proof}[Proof of Theorem~\ref{thm:polygon-einstein}]
%\label{sec:einstein-derivation}% label on proof block; use thm ref instead
The proof has four steps: (1)~force the interaction to be
logarithmic, (2)~extract $R(x)$ from the Hadamard parametrix,
(3)~conclude $R = \mathrm{const}$ from (H2) $+$ Step~2,
(4)~promote to Einstein in any spacetime dimension.

\medskip
\noindent\emph{Step~1 (Logarithmic interaction from (H1)$+$(H3)).}
\textbf{Polynomial Casimir forces logarithmic interaction.}
Condition~(H1) requires the angular Casimir
to be the exact polynomial $f(m,N) = m(N{-}m)/2$.
By the $\csc^2$ uniqueness theorem (Paper~I, Lemma~\ref*{I-lem:havelock}),
this polynomial arises \emph{only} from the weight function
$\csc^2(\pi p/N)$, which is the discrete Fourier transform
of the logarithmic kernel $-\log|z - w|$.
Condition~(H3) (Hadamard universality) ensures that the Laplacian
Green's function has exactly this logarithmic singularity on any
Riemannian surface.
Together, (H1)$+$(H3) force the interaction potential to be
$V(z,w) = -G_S(z,w)$, where $G_S$ is the Green's function
of the scalar Laplacian on~$(S, g)$.

\medskip
\noindent\emph{Step~2 (Jensen maximisation).}
\textbf{Small-ring expansion reveals local curvature.}
The Green's function on any Riemannian surface has the short-distance
expansion $G_S(z,w) = -(1/2\pi)\log d(z,w) + \gamma_S(z) + O(d^2)$,
where $d(z,w)$ is the geodesic distance and $\gamma_S$ encodes
the global geometry.
Summing over the $N$-ring at radius~$\varepsilon$ centred at~$x$:
\begin{equation}\label{eq:C1-expansion-proof}
  C_1(x, \varepsilon)
  = (N{-}1)\!\left[1 - \frac{R(x)\,\varepsilon^2}{4}
    + O(\varepsilon^4)\right],
\end{equation}
where $R(x)$ is the scalar curvature at~$x$.
The mode-independence of the $O(\varepsilon^2)$ correction is
\emph{automatic}: $C_1$ is mode-independent by the Riemannian
Havelock identity (Paper~I, Theorem~\ref*{I-thm:riemannian-havelock}),
so its small-$\varepsilon$ expansion inherits this property
at every order.
The coefficient $1/4$ is proved from the exact $C_1$ formulas
on the three model surfaces.  On $\mathbf{H}^2$:
$C_1/(N{-}1) = (1{+}\xi^2)/(1{-}\xi)^2$ with
$\xi = \tanh^2(\varepsilon/2) \approx \varepsilon^2/4$,
giving $C_1/(N{-}1) \approx 1 + \varepsilon^2/2
= 1 - R\varepsilon^2/4$ (since $R = -2$).
On $\mathbf{S}^2$:
$C_1/(N{-}1) = (1{-}\xi)/(1{+}\xi)$ with
$\xi = \tan^2(\varepsilon/2) \approx \varepsilon^2/4$,
giving $1 - \varepsilon^2/2 = 1 - R\varepsilon^2/4$
(since $R = +2$).
On the plane: $C_1 = N{-}1$ (no correction; $R = 0$).
All three model surfaces give
$C_1/(N{-}1) = 1 - R\varepsilon^2/4 + O(\varepsilon^4)$
with the same coefficient $1/4$.
By the Riemannian Havelock identity
(Paper~I, Theorem~\ref*{I-thm:riemannian-havelock}),
$C_1$ is the unique surface-dependent coefficient in the
eigenvalue decomposition, so the $1/4$ extends to every
constant-curvature surface.

\medskip
\noindent\emph{Step~3 (Jensen maximisation $+$ Onsager selection).}
\textbf{(H2) $+$ Step~2 force $R = \mathrm{const}$.}
By the mathematical content of~(H2):
$S_{\mathrm{total}}[g] = \int S(x)\,dA$ is strictly
concave in~$C_1(x)$, so Jensen's inequality gives a
unique entropy maximum at uniform~$C_1$.
The Onsager principle (the physical content of~(H2))
selects this maximum: at negative temperature, the
partition function is dominated by the macrostate
maximising~$S_{\mathrm{total}}$.
Write $C_1^{\mathrm{eq}}$ for this uniform value.
Substituting the expansion~(\ref{eq:C1-expansion-proof}):
\[
  (N{-}1)\!\left[1 - \frac{R(x)\,\varepsilon^2}{4}
  + O(\varepsilon^4)\right] = C_1^{\mathrm{eq}}
  \quad\text{for all } x.
\]
The right side is $x$-independent.
Therefore $R(x) = R_0$ for all~$x$.
The \emph{value} of $C_1^{\mathrm{eq}}$ (and hence $R_0$ and
$\Lambda$) is determined separately by the energy constraint:
the WDW equation (Proposition~\ref{prop:wdw-lambda}) gives
$C_1^{\mathrm{eq}} = f(m^*,N)$ at the BO threshold,
with $R_0 = -4[f(m^*,N) - (N{-}1)]/[(N{-}1)\varepsilon_0^2]$.

\noindent
%\label{sec:vortex-gravity}% dead label, never referenced
Place a test $N$-gon at each point $x \in S$.
The frozen-mode entropy at~$x$ is
$S(x) = \tfrac{1}{2}\sum_{m \ne m^*}\log\lambda_m(x)$,
which depends on $R(x)$ through the Hadamard expansion
$C_1(x) = (N{-}1)[1 - R(x)\varepsilon^2/4 + O(\varepsilon^4)]$.
The entropy $S(x) = \tfrac{1}{2}\sum\log\lambda_m$ is
exact: the CMS angular modes
are integrable (Liouville--Arnold, \S\ref{sec:three-layer}),
so the microcanonical entropy equals the one-loop (log-determinant)
entropy---the higher-order corrections vanish identically because
the $N{-}2$ angle-action variables decouple.
$S$ is strictly concave in $C_1$:
$d^2S/dC_1^2 = -\tfrac{1}{2}\sum_{m \ne m^*}1/\lambda_m^2 < 0$.
The total entropy $S_{\mathrm{total}} = \int_S S(x)\,dA$
is therefore maximised (Jensen's inequality) when $C_1(x)$ is
uniform---i.e.\ when $R(x) = \mathrm{const}$.
The energy constraint $E = \int V(C_1(x))\,dA$ is compatible:
$V = \Phi(\rho) - f_{m^*}$ is a \emph{monotone}
function of $C_1$ (since $V' = \coth(\rho) > 0$), so fixing
$\bar{V} = E/\mathrm{Vol}$ is equivalent to fixing
$\bar{C}_1$ via the monotone map $V \mapsto C_1$.
Jensen's inequality on the concave functional $S(C_1)$
at fixed $\bar{C}_1$ gives the unique maximum at
$C_1(x) = \bar{C}_1$ for all~$x$.
The Lagrange multiplier for this energy constraint is~$\Lambda$.

\emph{The $O(\varepsilon^4)$ correction does not spoil the Jensen argument.}
The Hadamard expansion has the form
$C_1 = (N{-}1)[1 - R\varepsilon^2/4 + \alpha R^2\varepsilon^4 + O(\varepsilon^6)]$
with $\alpha = 1/48$ (the $a_1$ Seeley--DeWitt coefficient).
The second derivative of $S(C_1(R))$ with respect to~$R$ is
\[
  \frac{d^2S}{dR^2}
  = \underbrace{S''(C_1)(C_1')^2}_{\text{concave: $< 0$}}
  + \underbrace{S'(C_1)\,C_1''}_{\text{correction: sign$(\alpha)$}}.
\]
At leading order in~$\varepsilon^4$, strict concavity requires
$(N{-}1)\,\sigma_2/(32\alpha\,\sigma_1) > 1$,
where $\sigma_1 = \sum' 1/\lambda_m$ and
$\sigma_2 = \sum' 1/\lambda_m^2$ are spectral sums
at the threshold.
This ratio is computed for all $N = 7, \ldots, 24$
(see Code Availability):
it equals $7.5$ at $N = 7$ (the minimum) and grows
monotonically with~$N$.
The concavity term dominates the $O(\varepsilon^4)$ correction
by at least $7{:}1$; the Jensen argument holds throughout
the small-ring regime.

The metric is not varied by hand---it varies
\emph{because it is a vortex degree of freedom}
(Remark~\ref{rmk:logical-status} above).

\emph{Extension to smooth metrics via Regge calculus.}
The vortex-gravity isomorphism applies to conical deficits
(point masses), raising the question: how does (H2) constrain
a \emph{smooth} metric $R(x)$?
The bridge is Regge calculus \citep{Regge1961}:
any smooth 2D Riemannian manifold $(S, g)$ is
approximated by a simplicial complex $(S_h, g_h)$
in which curvature is concentrated at $V$ vertices,
each carrying a deficit angle
$\varepsilon_v \approx R(v)\,A_v$
(with $A_v$ the dual area).
The convergence $g_h \to g$ as $h \to 0$ is proved
\citep{CheegerMullerSchrader1984}; the Green's function
$G_h \to G$ in the $H^1$-norm (standard finite-element
convergence).
Each Regge vertex is a point vortex
(circulation $\kappa_v = \varepsilon_v/(2\pi)$,
by the isomorphism), so the Regge system is a
$V$-vortex configuration.
The vertex positions are \emph{simultaneously}
vortex positions and Regge metric degrees of freedom
(on a fixed triangulation, vertex positions determine
edge lengths, hence the simplicial metric).

\begin{lemma}[Onsager prerequisites for the Regge system]
\label{lem:onsager-regge}
Let $S$ be a compact surface and $\mathcal{T}$ a
triangulation with $V$ vertices whose deficit angles
are all of the same sign
(positive curvature: $\chi > 0$; or all negative:
$\chi < 0$ on a compact quotient of~$\mathbf{H}^2$).
The Regge polygon system satisfies the three prerequisites
of the Onsager variational principle:
\begin{enumerate}
\item \emph{Compact phase space.}
  The configuration space $S^V$ is compact ($S$ compact).
  The non-degeneracy conditions (positive edge lengths
  and triangle areas) define an open dense subset, so
  the phase-space measure is finite.
\item \emph{Energy bounded above.}
  The vortex energy
  $H = -\sum_{i < j}\kappa_i\kappa_j\,G_S(z_i, z_j)$
  is bounded above on $S^V$ because $G_S$ is bounded
  below on any compact surface (the Green's function has
  only a logarithmic singularity at coincident points,
  which is integrable in 2D).
\item \emph{Negative temperature.}
  By (1) and (2), the density of states
  $\Omega(E) = \mathrm{Vol}(\{z : H(z) = E\})$
  is a decreasing function of~$E$ for $E > E^*$
  (Onsager 1949; rigorous: Kiessling 1993, Theorem~2.1).
  The microcanonical temperature $\beta = d\ln\Omega/dE$
  is negative in this regime.
\end{enumerate}
The Onsager principle is a statement about which
macrostates dominate the microcanonical ensemble---it
does not require specific equations of motion.
The microcanonical measure on $\{H = E\}$ is the
Riemannian volume of the energy level set in~$S^V$,
which is well-defined whenever $H$ is smooth on a
compact manifold; no dynamics is needed to define it
\citep{Kiessling1993}.
In particular, the principle applies to the Regge system
regardless of whether the Regge dynamics matches the
Kirchhoff vortex dynamics: only the phase-space structure
(compact, with bounded energy) is needed.
\end{lemma}

\noindent
\emph{Self-referential energy and conformal invariance.}
The Regge Hamiltonian
$H = -\sum \kappa_i\kappa_j\,G_{g(z)}(z_i, z_j)$
appears self-referential: both the circulations~$\kappa_v$
and the Green's function~$G_{g(z)}$ depend on the vertex
positions through the metric.
On a \emph{flat} background ($\Lambda = 0$),
the conformal invariance of the 2D Green's function
(Paper~II, Proposition~\ref{II-prop:conformal-invariance})
gives $G = -(2\pi)^{-1}\log|z{-}w|$ in conformal
coordinates, \emph{independent of the deficit angles}.
The Regge energy is therefore exactly the Kirchhoff
Hamiltonian with fixed couplings and fixed
interaction kernel---the standard Onsager/CLMP/Kiessling
theorems apply without modification.
On a \emph{curved} background ($\Lambda \ne 0$), the
Hadamard expansion (Step~2 above) gives
$C_1(x, \varepsilon) = (N{-}1)[1 - R(x)\varepsilon^2/4
+ O(\varepsilon^4)]$.
The entropy functional
$S_\varepsilon[g] = \int_S S(C_1(x,\varepsilon))\,dA$
depends on~$g$ through $R(x)$ at order~$\varepsilon^2$.
This functional is \emph{strictly concave}
in~$R(x)$ for every $\varepsilon > 0$
(from $d^2S/dC_1^2 < 0$ and the affine dependence
of~$C_1$ on~$R$).
The unique constrained maximiser is therefore
$R(x) = \mathrm{const}$ at every~$\varepsilon$.
Taking $\varepsilon \to 0$: the Hadamard expansion
becomes exact, and the maximiser of
$S_\varepsilon$ converges to the maximiser of
$S_0[g] = \int S((N{-}1)[1 - R\varepsilon^2/4])\,dA$
by $\Gamma$-convergence
(\citealt{DalMaso1993}, Proposition~6.8):
the integrands $S(C_1(x,\varepsilon))$ converge uniformly
in~$g$ as $\varepsilon \to 0$ (Hadamard), the domain
$\{g : \mathrm{Area}(S,g) = A\}$ is compact,
and the strict concavity of~$S$ in~$C_1$
gives uniqueness at each~$\varepsilon$, upgrading
subsequential convergence to full convergence.
The self-referential correction (the metric dependence
of~$G_S$) enters at $O(\varepsilon^2)$ and does not
affect the maximiser because the concavity-based
conclusion $R = \mathrm{const}$ holds at \emph{every}
order in~$\varepsilon$, not merely at leading order.
The Onsager prerequisites (compact phase space, bounded
energy, decreasing density of states at high energy)
are preserved throughout because the
$O(\varepsilon^2)$ perturbation is smooth and bounded
on the compact domain~$S$.

\noindent
Given Lemma~\ref{lem:onsager-regge}, the Onsager
principle on the $V$-vertex Regge system selects the
entropy-maximising vertex configuration, which by~(H2a)
has $C_1(v)$ uniform---hence $R = \mathrm{const}$.
In the continuum limit $V \to \infty$:
$S_{\mathrm{Regge}} \to S_{\mathrm{total}}[g]$
(by the FEM convergence of the Green's function,
sharpened to $O(h^\infty)$ by the Szeg\H{o}
estimates below),
and the entropy-maximising Regge metric converges to
the entropy-maximising smooth metric
(since $S_{\mathrm{poly}}$ is strictly concave with a
unique maximum, pointwise convergence of
$S_h \to S$ implies convergence of the
maximiser \citep{DalMaso1993}).
No new physical assumption enters: the Regge
approximation is a mathematical theorem, the
vortex-gravity isomorphism is kinematic, and
the Onsager prerequisites are verified in
Lemma~\ref{lem:onsager-regge}.

The Szeg\H{o}--Toeplitz convergence analysis
(supplementary material) sharpens the
FEM convergence rate to $O(h^\infty)$ using the
analyticity of the $\csc^2$ symbol.


%% Szegő-Toeplitz convergence analysis (184 lines) moved to
%% supplementary material. The analyticity of the csc² symbol
%% gives O(h^∞) Regge entropy convergence, sharpening the
%% O(h) FEM estimate. See the supplement for the full
%% Wiener-Hopf factorization and shell-dominance bounds.


\emph{Partition function interpretation.}
The total polygon partition function
$Z[g] = \sum_N e^{-b(N)}\,Z_{\mathrm{frozen}}(N, g)$
is a functional of the metric~$g$ through
$Z_{\mathrm{frozen}} = \prod_m |\lambda_m(g)|^{-1/2}$
(Proposition~\ref{prop:gravity-partition}).
By the vortex-gravity isomorphism, this is the
\emph{gravitational} one-loop partition function summed
over polygon sectors: $Z[g] = Z_{\mathrm{grav}}[g]$,
with $e^{-b(N)}$ playing the role of $e^{-S_{\mathrm{EH}}}$
(the classical Euclidean action of the $N$-cone geometry).
The saddle-point condition $\delta(-\log Z)/\delta g = 0$ is
therefore the Einstein equation---specifically, it is the
Euler--Lagrange equation for the combined classical action
$b(N)$ plus the one-loop correction $-\log Z_{\mathrm{frozen}}$.
This identification gives the WDW equation
(\S\ref{sec:wdw}) its physical meaning: it is the
\emph{quantum mechanics of the gravitational partition function}.
The WKB tunnelling action $S_{\mathrm{BO}}(N)$ is the
Euclidean action of the dominant instanton connecting the
two turning points of $V(\rho)$, and $\exp(-S_{\mathrm{BO}})$
is the tunnelling amplitude between the Planck-scale and
cosmological-scale vacua.
The hierarchy $\ell \cdot v = \exp(S_{\mathrm{BO}})$
(Paper~IV, \S14) is therefore not an \emph{ad hoc}
identification but a consequence of the partition function
interpretation: the cosmological scale $\ell$ is set by
the instanton suppression of the polygon partition function.

\emph{The explicit $\delta S/\delta g$ computation.}
In 2D conformal gauge $g = e^{2\varphi}\,g_0$, the scalar
curvature satisfies $\delta R/\delta\varphi = -2R$
and the area element $\delta(dA)/\delta\varphi = 2\,dA$.
Using the chain $dS/dC_1 = (1/2)\sum 1/\lambda_m$ (the
entropy gradient) and $dC_1/dR = -(N{-}1)\varepsilon^2/4$
(the Hadamard expansion, Step~2), the functional
derivative of $S_{\mathrm{total}}$ is:
\[
  \frac{\delta S_{\mathrm{total}}}{\delta\varphi(x)}
  = \frac{(N{-}1)\varepsilon^2}{4}\,R(x)\sum\frac{1}{\lambda_m(x)}
  + 2\,S(x).
\]
At the entropy maximum, this is constant for all~$x$:
the Euler--Lagrange equation is $F(R(x)) = \Lambda_\varphi$
where $F(R) = [(N{-}1)\varepsilon^2/4]\,R\sum(1/\lambda_m)
+ 2S$.
The function~$F$ is strictly monotone in~$R$
(since $dS/dC_1 > 0$ and $dC_1/dR < 0$, so
$dF/dR \ne 0$).
A monotone function equal to a constant forces
$R(x) = R_0 = \mathrm{const}$.
No Jensen inequality is needed---the Euler--Lagrange
equation is solvable and its unique solution is
constant curvature.

\medskip
\noindent\emph{Step~4 (Higher dimensions).}
\textbf{Mode-independence + $R = \mathrm{const}$
implies Einstein in any spacetime dimension~$d_{\mathrm{st}}$.}
In $d_{\mathrm{st}} = 3$ ($2{+}1$D): the Weyl tensor has
$d_{\mathrm{st}}(d_{\mathrm{st}}{+}1)(d_{\mathrm{st}}{+}2)(d_{\mathrm{st}}{-}3)/12 = 0$
components, so $R = \mathrm{const}$ alone gives
$R_{\mu\nu} = \Lambda g_{\mu\nu}$.

In $d_{\mathrm{st}} \ge 4$: $R = \mathrm{const}$ gives only the
\emph{trace} of the Einstein equations.
However, condition~(H1) provides the missing information.
The Hadamard parametrix in $d$ dimensions gives
$G(x,y) = U(x,y)/\sigma^{(d-2)/2} + V\log\sigma + W$,
where $\sigma = d(x,y)^2/2$.
The Weyl tensor $W_{ijkl}$ enters the angular structure of~$V$
at order~$\sigma$ through the traceless Riemann component,
producing an $\ell = 2$ (quadrupole) angular dependence.
In $3{+}1$D spacetime ($d_{\mathrm{st}} = 4$), the Weyl tensor has
$d_{\mathrm{st}}(d_{\mathrm{st}}{+}1)(d_{\mathrm{st}}{+}2)(d_{\mathrm{st}}{-}3)/12 = 10$
independent components;
the $\ell = 2$ sector contributes $2\ell + 1 = 5$ of these
(the traceless symmetric part).
The $\mathbb{Z}_N$ Fourier modes $m = 1, \ldots, N{-}1$
project onto specific angular harmonics.
Due to the palindromic symmetry
$\cos(2\pi pm/N) = \cos(2\pi(N{-}p)m/N)$,
the number of \emph{independent real} angular probes is
$\lfloor N/2\rfloor$, not $N{-}1$.
To constrain all $5$ quadrupole components of the
traceless Ricci tensor in spatial dimension $d = 3$
($d_{\mathrm{st}} = 4$), we need
$\lfloor N/2\rfloor \ge 5$, i.e.\ $N \ge 10$ (even) or
$N \ge 11$ (odd).
The $N \ge 10$ polygon is \emph{unstable} ($N > N_{\mathrm{crit}} = 7$),
but this does not affect the argument: condition~(H1) is a
statement about the \emph{structure} of the Green's function
(mode-independence of~$C_1$), not about the \emph{sign} of the
eigenvalues.
The Havelock identity $\lambda_m = C_1 - m(N{-}m)/2$ holds for
all~$N$, regardless of stability
(Paper~I, Lemma~\ref*{I-lem:havelock}).
The polygon is used as a \emph{mathematical probe} of the angular
structure, not as a physical equilibrium.
Verified numerically: the $5 \times 5$ cosine matrix
$M_{mp} = \cos(2\pi pm/N)$ has rank~$5$ for $N = 10$ and $N = 11$,
but rank~$3$ for $N = 6, 7$ (insufficient).
Any nonzero traceless Ricci component therefore makes
$C_1$ depend on~$m$ for $N \ge 10$, violating~(H1).
Hence (H1) at every~$x$ with $N \ge 10$ forces
$R_{ij} - (R/d)\,g_{ij} = 0$ on each spatial slice.
(In spatial dimension $d = 3$, the spatial Weyl tensor
vanishes identically, so the $5$ independent components
of the traceless part of the Riemann tensor are all
traceless Ricci --- and these are exactly what the
polygon's $\lfloor N/2 \rfloor$ angular modes probe.)
The promotion to the full spacetime Einstein equations
$R_{\mu\nu} = \Lambda\,g_{\mu\nu}$ (including time-space
components) requires additionally the boundary Virasoro
mechanism of Paper~IV, which supplies the dynamical
content in the time direction.

In $d_{\mathrm{st}} = 5$ ($4{+}1$D), additional quadrupole
components appear, requiring $N \ge 12$.
\end{proof}

\begin{remark}[The logical structure of~(H2)]
\label{rem:H2-structure}
Three points merit emphasis.
\begin{enumerate}
\item \emph{Integrable systems do not relax.}
  Correct: we do not claim that the CMS angular modes
  relax or thermalize.
  The Liouville--Arnold theorem gives
  equidistribution on invariant tori
  (stronger than relaxation: exact confinement).
  The entropy maximisation in~(H2) comes from the
  \emph{Onsager selection principle}
  (which macrostate dominates the microcanonical ensemble
  at negative temperature), not from relaxation dynamics.
\item \emph{Two variational principles.}
  Angular equidistribution at fixed geometry and metric
  selection by the Onsager principle are logically
  \emph{separate} claims, now stated separately as~(H2a)
  and~(H2b).
  Neither implies the other:
  (H2a) is a mathematical characterisation of the entropy
  functional;
  (H2b) is the physical principle that selects the
  entropy-maximising metric.
\item \emph{The Born--Oppenheimer bridge.}
  At fixed~$g$, the angular modes determine
  $S(x,g)$ at each point
  (the frozen-mode entropy from the BO partition function).
  Varying~$g$ varies $S_{\mathrm{poly}}[g] = \int S(x,g)\,dA$.
  The Onsager principle selects the~$g$ that maximises
  $S_{\mathrm{poly}}$.
  The BO decomposition
  (fast angular modes, slow breathing mode~$\rho$)
  provides the mechanism: the BO effective potential for~$\rho$
  is $V_{\mathrm{BO}}(\rho) = V_0(\rho) - S(\rho)$, so
  maximising~$S$ at fixed energy is equivalent to
  minimising $V_{\mathrm{BO}}$
  ---the standard BO equilibrium condition.
\end{enumerate}
\end{remark}

\begin{proposition}[The cosmological constant from the WDW equation]
\label{prop:wdw-lambda}
The value of $\Lambda$ in
Theorem~\ref{thm:polygon-einstein} is determined by the
Wheeler--DeWitt equation for the breathing mode~$\rho$:
\[
  \Bigl[-\frac{1}{2c}\,\frac{d^2}{d\rho^2}
  + V(\rho)\Bigr]\Psi = 0,
  \qquad V(\rho) = \Phi(\rho) - f(m^*,N).
\]
Two independent ingredients determine this equation:

\emph{The mechanism} (why $\dot\rho^2$ appears).
The Kirchhoff equations are first-order; the breathing
mode~$\rho$ is a collective coordinate with no canonical momentum.
The $N{-}2$ angular modes form a harmonic bath whose
influence functional generates an effective kinetic term
for~$\rho$ (Caldeira--Leggett mechanism).
The Onsager--Machlup noise kernel gives
$S_{\mathrm{OM}} = \int[c/2\;\dot\rho^2 + V]\,d\tau$.
The CL conditions hold: (i)~the angular modes are
harmonic (exact Fourier diagonalisation);
(ii)~the coupling is mode-independent
($\Phi' = \coth\rho$ for all~$m$);
(iii)~timescale separation
($\tau_{\mathrm{tunnel}} \sim e^{S_{\mathrm{BO}}} \gg 1$).
The bath has $N{-}2$ discrete modes; the CL derivation is
exact for any number of harmonic oscillators with linear coupling
(Caldeira \& Leggett 1983, eq.~3.1), not only in the
continuum limit.

\emph{The coefficient} (why $c = 12\,b(N)$).
The coefficient $c$ is not the CL friction
$\gamma = \coth^2\!\rho\,\Sigma(1/\lambda_m) \sim O(N)$;
it is the algebraic invariant of the eigenvalue decomposition.
The frozen-mode entropy is
$S_{\mathrm{bath}}(\rho)
= \tfrac{1}{2}\sum_{m \ne m^*}\log\lambda_m(\rho)$,
and $b(N)$ is the mode-independent offset in the
eigenvalue decomposition $\lambda_m = \Phi - f_m + \delta_m$
with $\Phi(\rho) = \log(2\sinh\rho) + b(N)$
(Proposition~\ref{prop:equiv-rr}).
The coefficient $c = 12\,b(N)$ is the unique value for which
the WDW ground-state energy reproduces the classical Havelock
spectrum: if $c$ were different, the quantum corrections to
$V(\rho)$ would break the mode-independence of~$\Phi$
(the defining property of the Havelock decomposition).
In the classical limit $c \to \infty$:
$V(\rho^*) = 0$, giving $\Phi(\rho^*) = f(m^*,N)$ and hence
$\Lambda = -4[f(m^*) - (N{-}1)]/[(N{-}1)\varepsilon^2]$.
\end{proposition}

\begin{corollary}[Sign of the cosmological constant from the polygon number]
\label{rmk:lambda-from-N}
The sign of $\Lambda$ is fixed by the polygon:
$N \le 6$: $f(m^*) < N{-}1$, so $\Lambda > 0$
(positive curvature, $S^2$-like---surface curvature destabilises);
$N = 7$: $f(3,7) = 6 = N{-}1$, so $\Lambda = 0$ (flat);
$N \ge 8$: $f(m^*) > N{-}1$, so $\Lambda < 0$
(negative curvature, $\mathbf{H}^2$-like---surface curvature stabilises).
The flat threshold $N = 7$---Havelock's 1931
stability boundary---sits exactly at $\Lambda = 0$.
\end{corollary}

\begin{remark}[Relation to Schur's lemma]
In $2{+}1$D, Step~4 of the proof is a special case
of Schur's lemma (isotropic sectional curvature implies
constant curvature).
The \emph{added value} of the Havelock derivation is fourfold:
(a)~Steps~1--4 are new: polygon stability implies isotropic
Green's function, which Schur alone cannot provide;
(b)~the curvature value $\Lambda$ is \emph{quantised}
by the polygon number~$N$ through the explicit formula
$\Lambda = -4[f(m^*,N) - (N{-}1)]/[(N{-}1)\varepsilon^2]$;
(c)~the hypothesis of Schur's lemma (isotropy) is
\emph{derived} from the Hadamard universality of the
Green's function, not assumed;
(d)~in $d \ge 4$, the Havelock condition (H1) provides a
testable criterion (mode-independence of~$C_1$) equivalent
to the full Weyl-vanishing condition of Schur's lemma,
connecting the abstract curvature condition to
the observable stability threshold.
\end{remark}

\begin{remark}[Characterisation theorem in $d \ge 3$]
\label{rmk:higher-d}
Steps~1--3 are dimension-independent:
the $\csc^2$ uniqueness forces a logarithmic interaction
in any~$d$ (Step~1); the Hadamard parametrix
$G(x,y) = U/\sigma^{(d-2)/2} + V\log\sigma + W$ contains
$R(x)$ at order~$\varepsilon^2$ in every~$d$ (Step~2);
and (H3) forces $R = \mathrm{const}$ (Step~3).
Step~4 closes differently by dimension:

\medskip
\begin{center}
\small
\begin{tabular}{lcl}
\toprule
Dimension & Weyl components & Status \\
\midrule
$d = 2$ ($2{+}1$D) & $0$ &
  \textbf{Theorem}: (H1) proved by the Havelock identity \\
$d = 3$ ($3{+}1$D) & $10$ &
  \textbf{Spatial equivalence}: (H1) constrains traceless Ricci;
  full spacetime requires Paper~IV \\
$d \ge 4$ & $\ge 35$ &
  Same, with more traceless Ricci components \\
\bottomrule
\end{tabular}
\end{center}
\medskip

\noindent
In $d = 2$: the Havelock identity \emph{proves} (H1)
analytically---no assumption is needed.
In $d \ge 3$: condition~(H1) (mode-independence of~$C_1$) is
\emph{equivalent} to the spatial Green's function being isotropic,
which is equivalent to maximal symmetry, which is equivalent
to the vacuum Einstein equations with~$\Lambda$.
The soft Havelock decomposition
(Proposition~\ref{prop:soft-havelock})
shows that (H1) is \emph{approximately} satisfied on any background,
with residual $|R_m| < a|f_m|$ (verified for $N = 5,\ldots,20$
and $d = 3, 4, 5$):
the violation of~(H1) is controlled and measures
the departure from Einstein geometry.
The extension to full Riemann curvature (beyond Ricci) requires higher-rank tensor probes and is beyond the scope of this paper.
\end{remark}

\paragraph{Supporting evidence: spectral focusing.}
By Theorem~\ref{thm:rg-monotonicity} and
Corollary~\ref{cor:c-theorem},
$N_{\mathrm{crit}}(\Delta)$ is non-increasing:
once a polygon type becomes unstable,
it remains unstable for all larger~$\Delta$.

\noindent
The monotonicity is formally analogous to
Raychaudhuri focusing (a parallel in structure, not a theorem):
\begin{center}
\small
\begin{tabular}{lll}
\toprule
Concept & General relativity & Polygon system \\
\midrule
Flow parameter & Affine parameter $\lambda$ & Conformal dim $\Delta$ \\
Expansion & $\theta = \nabla \cdot k$ (decreasing) &
  $N_{\mathrm{crit}}(\Delta)$ (decreasing) \\
Focusing & $d\theta/d\lambda \le 0$ &
  $dN_{\mathrm{crit}}/d\Delta \le 0$ (monotone) \\
Trapped & $\theta < 0$ (trapped surface) &
  $\lambda_m < 0$ (unstable mode) \\
Energy condition & $R_{\mu\nu}k^\mu k^\nu \ge 0$ &
  $C_1(\xi) > 0$ (spectral gap) \\
\bottomrule
\end{tabular}
\end{center}

\paragraph{The palindromic threshold as a Killing horizon.}
The WDW equation
$[-1/(2c)\,\Psi'' + V(\rho)\,\Psi = 0]$
is the zero-energy Schr\"odinger equation in the potential~$V$.
The turning point $V(\rho^*) = 0$ is a \emph{Killing horizon}
in the minisuperspace geometry: the Euclidean continuation
around~$\rho^*$ requires periodicity
$\beta = 2\pi/|V'(\rho^*)| = 2\pi/\coth(\rho^*)$
to avoid a conical singularity
\citep{GibbonsHawking1977}.
This is the same calculation that gives the de~Sitter temperature
$T = H/(2\pi)$---with the breathing mode~$\rho$ playing the
role of the scale factor~$a$.
The Gibbons--Hawking temperature at the palindromic threshold is
\begin{equation}\label{eq:T-GH}
  T = \frac{V'(\rho^*)}{2\pi} = \frac{\coth(\rho^*)}{2\pi}.
\end{equation}
\emph{Logical role of the Clausius relation.}
The vacuum Einstein equations $R = \mathrm{const}$ are derived
\emph{independently} from entropy maximisation
(Theorem~\ref{thm:polygon-einstein}, Steps~1--3).
The Clausius relation at the Gibbons--Hawking horizon
serves a \emph{different} purpose: it extracts the
gravitational coupling~$G_{\mathrm{eff}}$ from the
constant-curvature geometry.
On a constant-curvature surface, there is one threshold
(one~$N$) and one temperature~$T$.
The Jacobson argument \citep{Jacobson1995,Lashkari2014}
then promotes $G_{\mathrm{eff}}$ to the matter coupling:
$\delta Q = T\,\delta S$ at the Killing horizon gives
the full Einstein equations $G_{\mu\nu} + \Lambda g_{\mu\nu}
= 8\pi G_{\mathrm{eff}}\,T_{\mu\nu}$.

The three thermodynamic quantities are determined by
the BO potential $V(\rho) = \Phi(\rho) - f(m^*,N)$:
\begin{itemize}
\item $T = V'(\rho^*)/(2\pi) = \coth(\rho^*)/(2\pi)$
  (Unruh temperature from the surface gravity $\kappa = V'$);
\item $\delta Q = \kappa \times A \times \delta\rho
  = \coth(\rho^*) \times 2\pi\sinh(\rho^*) \times \delta\rho$
  (surface gravity $\times$ horizon circumference $\times$ displacement);
\item $\delta S = (dS/d\rho) \times A \times \delta\rho$
  where $A = 2\pi\sinh(\rho^*)$ is the horizon circumference and
  $dS/d\rho$ is the Bekenstein--Hawking entropy production rate.
\end{itemize}

\emph{Which entropy?}
The one-loop entropy $S' = -\tfrac{1}{2}\log|\det{}'H|$ (reduced
determinant, zero modes excluded) has a well-defined derivative
at~$\rho^*$:
\begin{equation}\label{eq:dS-spectral}
  \frac{dS'}{d\rho}\bigg|_{\rho^*}
  = -\frac{\coth(\rho^*)}{2}\,\sigma(N),
  \qquad
  \sigma(N) = \!\sum_{\substack{m=1 \\ f_m \ne f_{m^*}}}^{N-1}
  \frac{1}{f_{m^*} - f_m},
\end{equation}
where the sum runs over non-zero modes
($\sigma(7) = 8/3$, $\sigma(8) = 49/9$, verified $N = 7,\ldots,15$;
see Code Availability).
The zero mode's crossing contributes to the spectral flow
(APS eta-invariant shift $= \pm 1$ per palindromic pair),
not to $|{\det{}'H}|$.

\emph{Non-trivial Clausius relation via Cardy $+$ RT.}
The Bekenstein--Hawking entropy $S_{\mathrm{BH}} = A/(4G)$
presupposes~$G$, so using it in the Clausius relation gives
a tautology.
A non-trivial determination uses the \emph{Cardy thermal entropy}
of the boundary CFT, which depends on~$c$ but not on~$G$:
\begin{equation}\label{eq:cardy-entropy}
  S_{\mathrm{Cardy}}
  = \frac{\pi c}{3}\,T\,L
  = \frac{\pi c}{3}\,\cosh(\rho^*),
\end{equation}
where $T = \coth(\rho^*)/(2\pi)$ is the Unruh temperature,
$L = 2\pi\sinh(\rho^*)$ is the ring circumference,
and $c = 12\,b(N)$ is computed from the spectrum (Paper~I).
No free parameter enters.
The Ryu--Takayanagi formula $S = A/(4G)$ with
$A = 2\pi\sinh(\rho^*)$ then determines~$G$:
\begin{equation}\label{eq:clausius}
  G_{\mathrm{RT}}
  = \frac{A}{4\,S_{\mathrm{Cardy}}}
  = \frac{3\,\tanh(\rho^*)}{2\,c}
  = G_{\mathrm{BH}} \times \tanh(\rho^*).
\end{equation}
The correction $\tanh(\rho^*) = 1 - 2e^{-2\rho^*} + O(e^{-4\rho^*})$
is the standard BTZ finite-horizon correction, vanishing
exponentially: $6\%$ at $N = 7$, $0.03\%$ at $N = 11$,
$< 10^{-7}$ for $N \ge 15$
(verified numerically for $N = 7,\ldots,20$; see Code Availability).
The common factor $\coth(\rho^*) = V'(\rho^*)$ is not a
free choice: it is determined by the $\csc^2$ kernel
($V = \log(2\sinh\rho)$ is the $\mathbf{H}^2$ Green's function).
The one-loop spectral sum~\eqref{eq:dS-spectral} provides
an $N$-dependent quantum correction beyond the Cardy entropy;
incorporating it would improve the finite-$\rho^*$ convergence.
The \citet{Lashkari2014}
entanglement argument shows that positivity of relative entropy
$S(\rho\|\sigma) \ge 0$ in the unitary boundary CFT implies
$\delta Q = T\,\delta S$ at every Rindler horizon.
The Clausius relation is verified numerically at all
palindromic thresholds $N = 7, \ldots, 17$
(see Code Availability).
The extension to the full matter-coupled Einstein equations
$G_{\mu\nu} + \Lambda g_{\mu\nu} = 8\pi G\,T_{\mu\nu}$
requires the Jacobson thermodynamic argument
(\citealt{Jacobson1995}; \citealt{Lashkari2014}):
$\delta Q = T\,\delta S$ at every local Rindler horizon
implies Einstein with matter.
The identification of palindromic thresholds as Rindler
horizons uses the modular Hamiltonian of the angular bath,
which equals the Rindler energy by the
Bisognano--Wichmann theorem for the boundary CFT
(Paper~IV, Theorem~\ref*{IV-thm:nonlinear-einstein}).

\begin{proposition}[Partition function verification via the Gauss product]
\label{prop:partition-verification}
The polygon partition function
$Z_{\mathrm{poly}} = \sum_N e^{-b(N)}\,Z_{\mathrm{frozen}}(N)$
equals the Chern--Simons partition function
$Z_{\mathrm{CS}}$ on the Seifert manifold, with three
independent mathematical inputs:
\begin{enumerate}
\item \emph{Gauss product (algebra, Paper~I):}
  $b(N) = N(N{+}1)/12 + \frac{1}{N{-}1}\sum_{m=1}^{N-1}
  \ln\sin(\pi m/N)$.
  The transcendental part of~$b(N)$ is the
  \emph{mean log-sine} over the $N{-}1$ pairwise chord
  distances---the regularised self-energy from the identity
  $\prod_{m=1}^{N-1}\sin(\pi m/N) = N/2^{N-1}$.
\item \emph{Cheeger--M\"uller theorem (differential geometry):}
  the analytic torsion $\tau_{\mathrm{an}} = \prod|{\lambda_m}|$
  (spectral determinant of the gauge-covariant Laplacian)
  equals the Reidemeister torsion $\tau_R$ (combinatorial
  invariant of the Seifert manifold).
  For the $N$-gon flat connection:
  $Z_{\mathrm{frozen}}(N) = \tau_{\mathrm{an}}^{-1/2}$
  is the standard one-loop normalization.
\item \emph{CMS--CS Casimir identification
  (Proposition~\ref{prop:casimir-havelock}):}
  the vortex Hessian eigenvalues $\lambda_m = f(m^*){-}f(m)$
  are the eigenvalues of the cone Laplacian at the
  $\mathbb{Z}_N$-symmetric point (proved via the
  Helgason symmetric space theorem, \S\ref{sec:graviton}).
\end{enumerate}
The same Gauss product $\prod\sin(\pi m/N) = N/2^{N-1}$
appears on both sides:
\begin{itemize}
\item \emph{Vortex side}: it regularises the self-energy,
  giving the transcendental part of~$b(N)$;
\item \emph{Spectral side}: it determines the cone spectral
  zeta function $\zeta'_{\mathrm{cone}}(0)$, giving the
  analytic torsion.
\end{itemize}
They agree because both compute the regularised determinant
of the same $\mathbb{Z}_N$-circulant $\csc^2$ matrix---the
vortex side via the Havelock decomposition (Paper~I),
the spectral side via the zeta function
(Cheeger--M\"uller).
Points~(1) and~(2) are mathematical theorems.
Point~(3) is the only step using a physical identification
(the CMS--CS correspondence).
\end{proposition}

\paragraph{Matter coupling.}
On a quotient $\Gamma\backslash\mathbf{H}^2$, the Selberg
correction $\delta C_1$ shifts the threshold
$\rho^* \to \rho^*_{\mathrm{matter}}$, but the Clausius ratio
retains the form $4\pi\,\tanh(\rho^*_{\mathrm{matter}})$
(verified on the Bolza surface with $|\delta C_1| \le 0.021$).
The ratio~$4\pi$ is a local geometric identity (the metric is
locally~$\mathbf{H}^2$), giving
$G_{\mathrm{eff}} = 1/(16\pi)$ in natural units
($8\pi G_{\mathrm{eff}} = 1/2$ as $\rho^* \to \infty$).


%% ====================================================================
\section{The Onsager--Wheeler--DeWitt correspondence}
\label{sec:qm-gr-bridge}

The classical vortex system admits a natural
quantum-mechanical description through the Feynman--Kac
theorem, which identifies the Euclidean path integral
of the BO effective action with the WDW equation.
This is a \emph{mathematical identity} (not a physical
derivation of quantum mechanics from classical physics):
it shows that the polygon system's statistical mechanics
\emph{is} a quantum system, in the same sense that any
Euclidean path integral defines a transfer-matrix Hamiltonian.

\paragraph{The Feynman--Kac correspondence.}
The mapping has five steps, each either proved or a theorem:
\begin{enumerate}
\item \emph{Havelock diagonalisation (Paper~I, proved).}
  The angular Hessian of the $N$-gon is a circulant;
  Fourier diagonalisation gives $N{-}2$ \emph{independent
  harmonic oscillators} with frequencies $\lambda_m$.
\item \emph{Gaussian path integral (exact for harmonic oscillators).}
  Integration over the angular modes $\{a_m\}$ is
  an exact Gaussian: $\int\prod da_m\,
  \exp(-\tfrac{1}{2}\sum\lambda_m|a_m|^2)
  = \prod\lambda_m^{-1/2}$.
  No approximation is involved for a quadratic Hamiltonian.
\item \emph{Markovian effective action (timescale separation).}
  The influence functional from the bath decays on timescale
  $\tau_{\mathrm{bath}} \sim 1/\omega_m \sim O(1)$, while $\rho$
  tunnels on $\tau_{\mathrm{tunnel}} \sim e^{S_{\mathrm{BO}}}$.
  The local (Onsager--Machlup) approximation
  $S_{\mathrm{eff}} = \int[c/2\,\dot\rho^2 + V]\,d\tau$
  is valid to $O(e^{-S_{\mathrm{BO}}})$.
  The coefficient $c = 12\,b(N)$ is determined by the
  \emph{classical} Toeplitz positivity of the $\csc^2$ kernel
  (see below); no quantum input enters at this step.
\item \emph{Feynman--Kac theorem (1948, mathematical identity).}
  $Z = \int\mathcal{D}[\rho]\,e^{-S_{\mathrm{eff}}}$ has
  transfer matrix $T = e^{-\varepsilon H}$;
  $H\Psi_0 = 0$ is the WDW equation.
\item \emph{Born rule (consequence of path integral measure).}
  $|\Psi_0(\rho)|^2 = e^{-\beta V}/Z$ (the Boltzmann weight).
  In the WDW framework ($H\Psi = 0$ selects the unique
  physical state), this is the complete Born rule.
\end{enumerate}
The only \emph{physical} input is Step~1 (the Havelock
diagonalisation); the coefficient $c = 12\,b(N)$ in Step~3
is a \emph{classical} quantity (Toeplitz positivity of the
$\csc^2$ kernel, not a quantum computation).

\begin{proposition}[Self-decoherence of the breathing mode]
\label{prop:decoherence}
The $N{-}2$ angular Havelock modes decohere the
breathing mode~$\rho$ with decoherence rate
\begin{equation}\label{eq:decoherence-rate}
  \gamma = \frac{[\Phi'(\rho)]^2}{4}
  \sum_{m \ne m^*} \frac{1}{\lambda_m^2}.
\end{equation}
The reduced density matrix after tracing out the bath is
\[
  \rho_{\mathrm{red}}(\rho, \rho')
  = \Psi_0(\rho)\,\Psi_0^*(\rho')\,e^{-\gamma(\rho-\rho')^2}.
\]
For $\gamma \gg 1$ the off-diagonal elements are
exponentially suppressed and $\rho_{\mathrm{red}}$
is effectively diagonal:
$\rho_{\mathrm{red}} \approx |\Psi_0(\rho)|^2\,
\delta(\rho - \rho')$.
\end{proposition}

\begin{proof}
The CL influence functional for the Gaussian bath
with $\rho$-dependent frequencies $\omega_m(\rho)
= \sqrt{\lambda_m(\rho)}$ gives a decoherence kernel
$\exp(-\gamma(\rho{-}\rho')^2)$ with
$\gamma = \sum_m (g_m/\omega_m)^2$, where
$g_m = \partial\omega_m/\partial\rho
= \Phi'/(2\sqrt{\lambda_m})$ is the system-bath coupling.
Substituting: $\gamma = (\Phi')^2/4 \times
\sum 1/\lambda_m^2$.
For $N = 7$ (all modes stable on $\mathbf{H}^2$):
$\gamma \approx 7$--$12$ across the stable range.
For $N = 11$ at the stability threshold
$\rho^*_{11} \approx 0.88$ (the cosmologically relevant
operating point, where all $N{-}2 = 9$ bath modes
have $\lambda_m > 0$): $\gamma \gg 1$.
Near the threshold, the palindromic partner mode
($m = N{-}m^*$, which shares eigenvalue $\lambda_{m^*}$)
has $\lambda \to 0$, and $1/\lambda^2 \to \infty$;
this mode dominates the sum and
$\gamma$ grows without bound as the threshold is
approached.
The polygon is therefore \emph{strongly decohered}
at the physical operating point---not by fine-tuning
but as a structural consequence of the soft bath mode
near the stability threshold.

\emph{Irreversibility from near-integrability.}
The CMS angular bath is integrable \emph{in isolation}.
The coupling to the breathing mode~$\rho$ through the
$\rho$-dependent frequencies $\lambda_m(\rho)$
breaks exact integrability at $O(1/c)$
(the $W_N$ quartic coupling,
$\sim\!2\%$ at $N = 7$).
This converts exact Poincar\'e recurrences
(which would periodically restore coherence)
into diffusive mixing.
The Nekhoroshev theorem bounds the mixing
timescale from below:
$\tau_{\mathrm{mix}} \ge
\exp(c_0\,c^{a})$ for some $c_0 > 0$, $a > 0$
(with $\varepsilon = O(1/c)$ as the perturbation
parameter).
Even this exponentially long mixing time satisfies
$\tau_{\mathrm{mix}} \ll \tau_{\mathrm{tunnel}}
\sim e^{S_{\mathrm{BO}}}$,
provided $c_0\,c^a < S_{\mathrm{BO}}$
(for $c \approx 50$--$130$ and
$S_{\mathrm{BO}} \approx 38$--$103$, this is
satisfied for generic Nekhoroshev exponents
$a \le 1/2$).
Additionally, the Poincar\'e recurrence time
for the full $(N{-}2)$-mode system scales as
$\tau_{\mathrm{rec}} \sim \exp(O(N{-}2))$,
giving $\tau_{\mathrm{rec}} \sim e^5$--$e^9
\approx 150$--$8000$, also $\ll \tau_{\mathrm{tunnel}}$.
The decoherence is therefore \emph{irreversible}
on all cosmologically relevant timescales.

This is consistent with the exact BO result
(Paper~V, eq.~\eqref{V-eq:no-nonadiabatic}):
the mode-independence of~$\Phi$ makes the BO
potential exact (no non-adiabatic transitions),
while the $O(1/c)$ $W_N$ quartic---which breaks the
exact CMS integrability---provides the mode mixing
needed for irreversible decoherence.
The two results are \emph{compatible} because they
involve different aspects of the angular dynamics:
the BO exactness uses the Havelock structure
($\partial\lambda_m/\partial\rho$ mode-independent),
while the decoherence uses the nonlinear mode coupling
($W_N$ quartic at $O(1/c)$).
Neither contradicts the other.
\end{proof}

\noindent
\emph{Implications.}
The CL decoherence selects the $\rho$-basis as the pointer basis
($\gamma \gg 1$ at the physical operating point),
with $|\Psi_0(\rho)|^2 = e^{-\beta V}/Z$
as the Boltzmann weight \citep{Zurek2003}.
The polygon decoheres \emph{itself}: no external
observer or environment is required.
The $N{-}2$ angular modes form a Caldeira--Leggett bath for
the breathing mode~$\rho$.
The four CL conditions hold:
(i)~harmonic bath (Havelock diagonalisation, exact);
(ii)~uniform parametric coupling
($\partial\lambda_m/\partial\rho = \Phi'(\rho)$ for all~$m$);
(iii)~timescale separation
($\tau_{\mathrm{tunnel}} \sim e^{S_{\mathrm{BO}}} \gg \tau_m \sim O(1)$);
(iv)~FDT (automatic for Gaussian systems).
The resulting Onsager--Machlup action is
$S_{\mathrm{OM}} = \int[\tfrac{c}{2}\dot\rho^2 + V(\rho)]\,d\tau$,
where $\dot\rho^2$ is the noise kernel (not a kinetic energy).

The Markovian approximation
($\tau_{\mathrm{bath}} \sim O(1) \ll \tau_{\mathrm{tunnel}} \sim e^{S_{\mathrm{BO}}}$)
ensures corrections are $O(e^{-S_{\mathrm{BO}}})$;
the coefficient $c = 12\,b(N)$ is classical
(Proposition~\ref{prop:equiv-rr}).
The Feynman--Kac transfer matrix gives $H\Psi_0 = 0$
(the WDW equation, Step~4 above), with Born rule
$|\Psi_0|^2 = e^{-\beta V}/Z$ as a mathematical consequence:
\begin{center}
\small
\begin{tabular}{ll}
\toprule
Classical (Onsager) & Quantum (WDW) \\
\midrule
Boltzmann weight $e^{-\beta V}$ & $|\Psi_0|^2$ \\
Partition function $Z$ & $\langle\Psi_0|\Psi_0\rangle$ \\
Thermal correlation $\langle\rho(\tau_1)\rho(\tau_2)\rangle$
  & $\langle\Psi_0|\hat\rho\,e^{-H\Delta\tau}\hat\rho|\Psi_0\rangle$ \\
\bottomrule
\end{tabular}
\end{center}
\noindent
\begin{theorem}[Quantum structure from polygon dynamics]
\label{thm:qm-from-polygons}
Starting from the classical Kirchhoff equations for
$N$ point vortices on~$\mathbf{H}^2$, with no quantum
postulate, the following quantum-mechanical structures
are derived:
\begin{enumerate}
\item[\textup{(Q1)}] \emph{Hilbert space.}
  $\mathcal{H} = L^2(\mathbb{R},\,d\rho)$
  (Feynman--Kac, Step~4 above).
\item[\textup{(Q2)}] \emph{Hamiltonian.}
  $H = -\tfrac{1}{2c}\,d^2/d\rho^2 + V(\rho)$
  with $c = 12\,b(N)$ (CL mechanism $+$ orbifold matching).
\item[\textup{(Q3)}] \emph{Physical state.}
  The WDW constraint $H\Psi_0 = 0$ selects the ground
  state (the additive constant in~$V$ is fixed by
  $V(\rho^*) = 0$ at the BO equilibrium, giving
  $E_0 = 0$; excited states $E \ne 0$ do not satisfy
  the constraint).
\item[\textup{(Q4)}] \emph{Ground-state position distribution.}
  $|\Psi_0(\rho)|^2 = e^{-\beta V}/Z$
  (mathematical identity: the Boltzmann weight of the
  classical Onsager--Machlup path integral equals the
  quantum ground-state probability in the $\rho$-basis).
\item[\textup{(Q5)}] \emph{Pointer basis.}
  The angular bath decoheres~$\rho$ with
  $\gamma \gg 1$ near the stability threshold
  (Proposition~\ref{prop:decoherence}).
  The $\rho$-basis is the pointer basis.
\item[\textup{(Q6)}] \emph{Uncertainty principle.}
  $\Delta\rho\,\Delta p = 1/2$ (minimum-uncertainty
  ground state near the turning point;
  $\Delta\rho = 1/\!\sqrt{2c\omega}$,
  $\Delta p = \sqrt{c\omega/2}$,
  $\omega^2 = |V''(\rho_1)| = \mathrm{csch}^2(\rho_1)$).
\item[\textup{(Q7)}] \emph{Entanglement.}
  The entanglement entropy between~$\rho$ and
  the angular bath is $S_{\mathrm{ent}} \gg 1$~bit
  (from the CL squeezing parameters
  $r_m = (\Phi'/(2\lambda_m))^2$;
  see Code Availability).
\item[\textup{(Q8)}] \emph{Irreversibility.}
  The breathing-angular coupling breaks exact CMS
  integrability at $O(1/c)$, converting exact
  recurrences into diffusive mixing.
  The Nekhoroshev bound gives
  $\tau_{\mathrm{mix}} \ge \exp(c_0 c^a)$;
  even this exponentially long mixing time satisfies
  $\tau_{\mathrm{mix}} \ll \tau_{\mathrm{tunnel}}
  \sim e^{S_{\mathrm{BO}}}$ for generic exponents.
  Decoherence is irreversible on all cosmologically
  relevant timescales.
\end{enumerate}
Within the WDW framework ($H\Psi = 0$, unique state,
no external time), items \textup{(Q1)--(Q8)} constitute
the quantum-mechanical structure of the breathing mode:
no additional postulate (measurement, collapse, or
Born rule for arbitrary states) is needed for this
sector.
\end{theorem}

\begin{proof}
Items (Q1)--(Q4) follow from the five-step chain above
(Steps~1--5: Havelock $\to$ CL $\to$ Onsager--Machlup
$\to$ Feynman--Kac $\to$ WDW).
Item (Q5) is Proposition~\ref{prop:decoherence} with the
mode-mixing irreversibility.
Item (Q6): the WDW ground state near the turning point
$\rho_1$ is $\Psi_0 \propto
\exp(-c\omega(\rho{-}\rho_1)^2/2)$
with $\omega^2 = |V''(\rho_1)| = \mathrm{csch}^2(\rho_1)$;
the widths $\Delta\rho = 1/\!\sqrt{2c\omega}$ and
$\Delta p = \sqrt{c\omega/2}$ give
$\Delta\rho\,\Delta p = 1/2$ (minimum uncertainty).
Item (Q7): the CL influence functional at zero temperature
gives squeezing parameter
$r_m = (\Phi'/(2\lambda_m))^2$ for each bath mode;
the entanglement entropy is
$S_{\mathrm{ent}} = \sum_m [(r_m{+}1)\ln(r_m{+}1)
- r_m\ln r_m]$.
Item (Q8): the $O(1/c)$ $W_N$ quartic coupling
breaks exact CMS integrability; the Nekhoroshev
bound gives $\tau_{\mathrm{mix}} \ge \exp(c_0 c^a)$,
and $\exp(c_0 c^a) \ll e^{S_{\mathrm{BO}}}$
for generic exponents $a \le 1/2$ at
$c \approx 50$--$130$, $S_{\mathrm{BO}} \ge 38$.
\end{proof}

\paragraph{Summary of the bridge.}
\begin{center}
\small
\begin{tabular}{lcl}
\toprule
Output & Mechanism & Input \\
\midrule
Schr\"odinger (WDW) & Feynman--Kac & Onsager stat.\ mech.\ \\
CS Hilbert space & Vortex partition fn.\ at $k=c/6$ & $b(N)$ from eigenvalues \\
$3$D Einstein & Havelock + Hadamard & Paper~III, Thm.\ 5 \\
$4$D vacuum Einstein & CS + KK & Paper~IV, Theorem~\ref*{IV-thm:nonlinear-einstein}(a) \\
$4$D Einstein + matter & CS + KK + RT (derived from CFT$_2$) & Paper~IV, Theorem~\ref*{IV-thm:nonlinear-einstein}(b) \\
$\mathrm{SU}(3) \times \mathrm{SU}(2) \times \mathrm{U}(1)$ &
  McKay + Eightfold Way + KK & $N = 7$ \\
$\sin^2\theta_W = 3/11$ & CS conformal weights & Paper~IV, Theorem~\ref*{IV-thm:weinberg} \\
$\delta_{\mathrm{CKM}} = 70.2^\circ$ & BF-crossing phase & $\arg\Gamma(\tfrac{1}{2}+i)$ \\
\bottomrule
\end{tabular}
\end{center}
\noindent
Every row traces back to the classical vortex system.
The quantum-mechanical description of the breathing mode
enters through the Feynman--Kac theorem
(a mathematical identity relating Brownian motion to the
Schr\"odinger equation).


\paragraph{Discussion and predictions.}
The status classification, parameter accounting, cumulative
uncertainty analysis, and falsifiable predictions for the
results of this paper appear in Paper~VI.

\paragraph{Code availability.}
All numerical verifications and computational proofs are
available in the companion code repository.

\bibliographystyle{plainnat}
\bibliography{../shared/refs}
\end{document}
